\documentclass[12pt]{report}

\usepackage[margin=1in]{geometry}
\usepackage[T1]{fontenc}
\usepackage[utf8]{inputenc}
\usepackage{lmodern}
\usepackage{amsmath,amssymb}
\usepackage{graphicx}
\usepackage{booktabs}
\usepackage{longtable}
\usepackage{array}
\usepackage{tabularx}
\usepackage{caption}
\usepackage{hyperref}
\usepackage{setspace}
\IfFileExists{enumitem.sty}{%
    \usepackage{enumitem}
}{%
    \usepackage{xparse}
    % Fallback for lean TeX installs: preserve itemize syntax while ignoring enumitem options.
    \RenewDocumentEnvironment{itemize}{o}{%
        \begin{list}{\labelitemi}{%
            \setlength{\leftmargin}{1.8em}%
            \setlength{\itemsep}{0.25\baselineskip}%
            \setlength{\parsep}{0pt}%
            \setlength{\topsep}{0.25\baselineskip}%
        }%
    }{%
        \end{list}%
    }%
}

\hypersetup{colorlinks=true,linkcolor=black,citecolor=black,urlcolor=blue}
\graphicspath{{figures/}}
\setlength{\parindent}{0pt}
\setlength{\parskip}{0.6\baselineskip}
\setcounter{secnumdepth}{2}
\setcounter{tocdepth}{2}
\onehalfspacing

\newcolumntype{L}[1]{>{\raggedright\arraybackslash}p{#1}}
\newcolumntype{C}[1]{>{\centering\arraybackslash}p{#1}}
\newcolumntype{Y}{>{\raggedright\arraybackslash}X}

\newcommand{\corefigure}[4][0.88\textwidth]{%
    \begin{figure}[htbp]
        \centering
        \IfFileExists{figures/#2}{%
            \includegraphics[width=#1]{#2}%
        }{%
            \fbox{\parbox{0.82\textwidth}{\centering Placeholder for figure\\Expected file: \texttt{\detokenize{figures/#2}}\\Drop the file into the \texttt{figures/} folder to load it automatically.}}%
        }%
        \caption{#3}
        \label{#4}
    \end{figure}
}

\newcommand{\optionalfigure}[4][0.88\textwidth]{%
    \IfFileExists{figures/#2}{%
        \begin{figure}[htbp]
            \centering
            \includegraphics[width=#1]{#2}%
            \caption{#3}
            \label{#4}
        \end{figure}
    }{}%
}

\title{\textbf{For Whom Does Game-Based Sales Training Work?}\\
Expertise-Conditional Effects in an RCT of a Scaffolded B2B Simulation}
\author{Yoshua Yanottama}
\date{April 2026}

\begin{document}

\hypersetup{pageanchor=false}
\pagenumbering{roman}
\maketitle
\hypersetup{pageanchor=true}

\chapter*{Abstract}
\addcontentsline{toc}{chapter}{Abstract}
\textbf{Background:} Game-based learning (GBL) produces inconsistent outcomes in professional sales training, yet the moderating role of learner expertise remains underexplored.

\textbf{Methods:} A pre-registered randomised controlled trial ($N = 63$) compared a scaffolded B2B sales simulation to a conventional e-learning programme. Sales self-efficacy (SE) and adaptive selling beliefs (ADAPTS) were analysed using ANCOVA with pre-test scores as covariates, following CONSORT recommendations; cognitive load, engagement, and transfer intention were examined with Welch $t$-tests.

\textbf{Results:} ANCOVA revealed a marginal trend for SE ($F$(1,60) = 3.56, $p$ = .064, $\eta^2_p$ = .056, approaching medium effect) and a non-significant difference in ADAPTS ($F$(1,60) = 2.22, $p$ = .141). The game condition incurred substantially higher perceived cognitive load ($d$ = 0.71) without corresponding gains on the main post-test outcomes. Pre-registered continuous moderation showed that outcomes depended significantly on pre-intervention expertise (pre-ADAPTS $\times$ condition: $F$(1,59) = 11.65, $p$ = .0012; pre-SE $\times$ condition: $F$(1,59) = 5.52, $p$ = .022). An exploratory cluster analysis then illustrated this pattern: novice learners ($n = 10$) showed a large negative effect on adaptive selling gain ($d = -1.88$, $p < .001$), while intermediate and advanced clusters trended positively or neutrally. Within the game group, SE fully mediated the effect of engagement on adaptive selling (indirect effect = 1.10, 95\% CI [.43, 2.24]). Expertise moderation was corroborated by differential SE--adaptive selling correlations across conditions ($r = .877$ vs.\ $.532$, Fisher $z$, $p = .004$) and by behavioural telemetry (first-attempt latency: $r = -.37$, $p = .031$).

\textbf{Conclusions:} At the aggregate level, the game produced higher cognitive load without superior learning outcomes, constituting a null-to-marginal primary finding. Exploratory analyses provide preliminary evidence that effectiveness is expertise-conditional: novice learners are disadvantaged by the additional interface load imposed before domain schemas exist, while intermediate and advanced learners are neutral or marginally benefited. These exploratory findings require pre-registered replication, but support adaptive deployment strategies---such as prerequisite screening or dynamic difficulty adjustment---and extend expertise reversal theory to professional GBL contexts as a hypothesis warranting confirmatory investigation.

\textbf{Keywords:} digital game-based learning, cognitive load, randomised controlled trial, expertise reversal, sales self-efficacy, adaptive selling

\clearpage

\tableofcontents
\listoftables

\chapter*{Acknowledgements}
\addcontentsline{toc}{chapter}{Acknowledgements}
This manuscript consolidates the finalized statistical findings from the analysis report with the theoretical framing developed in earlier proposal documents. The document is written as a thesis-ready version aligned to the current evidence base.

\clearpage
\pagenumbering{arabic}
\chapter{Introduction}

\section{Background}
Digital game based learning (DGBL) has emerged as a promising design for learning interventions with better learning performance, improved attitude in learning, and higher motivation, in comparison to traditional face-to-face instruction (Chang et al., 2020) or even to online virtual learning environments (Hern\'andez and Moreno, 2019). It has been used in various fields and industries to improve critical thinking and argumentation skills (Noroozi et al., 2020), leadership skills (Sousa and Costa, 2014), and even as part of wider corporate competence development (Senderek et al., 2015). These benefits have led to its increasing acceptance in workplace training, yet barriers persist: lack of available solutions, lack of expertise in authoring, doubts about efficacy (Tay et al., 2022), and the negative social association of games in some settings (B\"osche and Kattner, 2011). As Tay et al. suggest, questions of efficacy need to be addressed through stronger alignment with pedagogical principles and teaching strategies, because without a coherent design framework DGBL can suffer from the inclusion of non-essential design and gamification components that potentially hinder the achievement of learning outcomes. Rooting DGBL in sound pedagogical and learning strategies through informed instructional design may improve its appeal in workplace training beyond mere novelty effect and strengthen the validity of claims made about its impact. A useful framework for this intentional design is the learning mechanics--game mechanics framework (LM--GM), which encourages alignment between game mechanics and intended learning outcomes through corresponding learning mechanics (Arnab et al., 2015).

There is growing interest in the potential of DGBL for corporate training, particularly within the demanding field of sales education (Cummins et al., 2013; Tay et al., 2022), but concerns about efficacy and the attribution of the factors leading to DGBL success cannot be answered adequately by many existing research designs, limiting the explanatory power of current findings and their utility in informing robust DGBL design (All et al., 2021). While many studies cite increased engagement as the key driver of improved learning outcomes over traditional methods (Dahalan et al., 2024; Woo, 2014), the specific factors facilitating that engagement have not been investigated in sufficient depth. Cognitive load theory (Sweller et al., 2011) and self-efficacy theory (Bandura, 1977, 1986) can be leveraged as promising explanatory factors for DGBL effectiveness in sales education and beyond (Barling and Beattie, 1983; Hung et al., 2014).

The DGBL evidence base also suffers from methodological limitations, because the scarcity of randomized controlled trials (RCTs) constrains reliable comparison with traditional designs and limits strong causal inferences regarding effectiveness in promoting learning outcomes (All et al., 2021; Hainey et al., 2016; Tay et al., 2022). All et al.'s (2021) systematic review of DGBL effectiveness assessment identifies several recurring design weaknesses that this study attempts to address: the failure of true random assignment in online recruitment contexts (motivating the use of alternating blocked allocation here), the insufficient attention to instructional time as a factor in attrition (directly evidenced in the present study's behavioral logs, where the majority of game dropouts occurred at phase transitions associated with the highest time demands), and the under-examined role of motivation to start playing. On the last point, All et al. note that entry motivation is a confound that is rarely measured or controlled; in the present study, it manifests precisely as selective attrition rather than as a baseline variable, allowing the dropout pattern itself to serve as behavioral evidence of the friction that low-entry-motivation participants encounter. These design choices and their consequences are discussed further in Section~\ref{sec:method} and Section~\ref{sec:h1}. Specifically, there is a need for more empirical investigation into how scaffolding strategies embedded within DGBL for sales training influence the interplay between cognitive load and self-efficacy (Chang et al., 2018; Evans et al., 2024), and how these factors collectively shape adaptive selling and related applied sales development (Cardinali et al., 2025). While traditional sales education methods face challenges in scalability and in providing safe, repeatable practice (Cummins et al., 2013), there remains a gap in research directly comparing a pedagogically grounded DGBL intervention with standard digital formats in the sales domain. The potential of specific DGBL modalities, such as digital card games (Kordaki and Gousiou, 2017), to address these issues in sales training also warrants closer research attention.

\section{Research Objective}
To guide and focus the investigation, the following research questions are formulated:

\begin{itemize}[leftmargin=1.2cm]
    \item \textbf{Research question 1 (RQ1):} How does a scaffolded game-based learning (DGBL) intervention compare to a traditional click-and-continue digital training method in terms of its impact on sales training outcomes?
\end{itemize}

This primary research question is explored by comparing the two training conditions across several key outcome variables. Specifically, the study examines whether the scaffolded DGBL intervention leads to significantly different levels of learner engagement, sales self-efficacy, perceived cognitive load, adaptive selling, and transfer intention. In addition, it evaluates whether the DGBL group demonstrates stronger improvement on adaptive-selling beliefs and related sales-training outcomes from pre-test to post-test relative to the traditional digital training group.

\begin{itemize}[leftmargin=1.2cm]
    \item \textbf{Research question 2 (RQ2):} Within the context of the scaffolded DGBL sales training, what are the relationships between perceived cognitive load, learner engagement, sales self-efficacy, and adaptive-selling performance?
\end{itemize}

This second research question focuses on the internal dynamics within the DGBL intervention group. It investigates the extent to which perceived cognitive load, learner engagement, and sales self-efficacy are interrelated. It also explores how these psychological variables, both individually and collectively, relate to adaptive selling, transfer intention, and broader sales-training outcomes. Finally, the study examines whether there is evidence that the influence of managed cognitive load on adaptive-selling and transfer outcomes is channeled through self-efficacy and/or learner engagement, thereby clarifying the potential mechanisms at work inside the DGBL environment.

\section{Significance of the Study}
This research is anticipated to make several contributions to the fields of educational technology, instructional design, corporate training, and sales education. Its significance lies in the potential to advance both theoretical understanding and practical application of Digital Game-Based Learning (DGBL).

\subsection{Theoretical Contributions}
The study provides an empirical test of a DGBL intervention designed explicitly around cognitive load theory (CLT), self-efficacy theory, motivation theory, and the structured LM--GM game design framework. The primary aggregate finding --- higher cognitive load in the game condition without corresponding gains in self-efficacy or adaptive selling --- constitutes a theoretically meaningful result in its own right: it documents a case in which a pedagogically aligned game design produced additional cognitive overhead rather than the expected efficiency benefits. The exploratory expertise-conditional pattern (novice disadvantage $d$ = $-$1.88, advanced learner neutrality-to-advantage) provides preliminary evidence that CLT's expertise reversal mechanism extends to professional DGBL contexts. These findings, taken together, offer a more nuanced account of when and for whom LM--GM-grounded DGBL can be expected to work, while identifying mechanics transparency as a design criterion the framework does not yet operationalise.

\subsection{Practical Contributions}
The study involves the creation and evaluation of a specific scaffolded, card-based digital game for sales training. The aggregate null-to-marginal finding is itself a practically informative result: it establishes that a well-designed, LM--GM-grounded game does not automatically outperform a conventional digital module for a mixed-expertise population. The exploratory expertise-conditional pattern suggests that deployment decisions should account for learner prior knowledge, with novice learners directed to direct instruction before game-based modules are introduced. The high attrition rate in the game condition (41.8\%) is an additional practical signal: organizations should anticipate non-trivial dropout rates in game-based online training and treat completion as a deployment outcome in its own right, not only as a statistical filter.

\subsection{Methodological Contributions}
The study makes several methodological contributions to GBL evaluation research. First, it demonstrates the feasibility of conducting a pre-registered randomised controlled trial of a DGBL intervention in a fully online, professional training context, pairing survey-based psychometric outcomes with behavioural telemetry from game logs --- a combination that is comparatively rare in the GBL literature (All et al., 2021). Second, the clustering approach used to operationalise expertise groups was deliberately restricted to pre-intervention features only, avoiding the circularity that arises when post-test or gain variables are included in cluster solutions; this methodological choice is foregrounded as a template for future expertise-stratified GBL analyses. Third, a continuous moderation analysis (pre-outcome $\times$ condition interaction regression) is presented as the primary evidence for expertise moderation, rather than relying solely on the post-hoc cluster solution; this approach avoids the discretisation assumptions inherent in $k$-means and is more directly replicable in future designs. Fourth, two intention-to-treat sensitivity analyses are reported alongside the per-protocol primary analysis, bounding the range of plausible treatment effects under different attrition assumptions. Fifth, the study contributes a methodological critique that extends beyond this study: immediate post-test comparison is structurally conservative for game-based conditions, because games operate on motivational and schema-building mechanisms whose outputs compound post-training rather than expressing fully at the moment of completion. Adopting delayed retention and transfer endpoints as primary --- rather than supplementary --- outcomes is proposed as a field-wide standard for comparative GBL research.

In summary, this research aims to bridge theory and practice by developing and evaluating a pedagogically grounded DGBL intervention. It is intended to provide knowledge that can improve the design and implementation of future DGBL initiatives, particularly in the critical area of corporate sales training.

\section{Delimitations of the Study}
This study's findings should be interpreted within the following scope:

\begin{itemize}[leftmargin=1.2cm]
    \item \textbf{Participant Generalizability:} While targeting novice sales professionals, open online recruitment means the sample may be diverse; results primarily apply to this self-selected group.
    \item \textbf{Intervention Specificity:} The study evaluates a specific scaffolded card-based DGBL (``Improve Your Sales Game''); findings may not generalize to all DGBL formats or all sales training contexts.
    \item \textbf{Outcome Scope:} The focus is on self-reported adaptive selling, transfer-oriented perceptions, and related immediate post-training outcomes, not the full spectrum of sales competencies or actual on-the-job performance.
    \item \textbf{Measurement Timeline:} Outcomes are assessed immediately post-intervention, without long-term follow-up on retention or real-world skill transfer.
\end{itemize}

\chapter{Literature Review}

\section{Sales Education and the Case for DGBL}

\subsection{Why DGBL Fits Sales Education}
The evolving landscape of business, the continuous demand for professional upskilling, and the characteristics of modern learners underscore the growing need for innovative pedagogical approaches such as DGBL in sales education. Sales roles are inherently performance-oriented and skills-based, requiring ongoing development to meet dynamic market demands (Cummins et al., 2013). DGBL offers a compelling response to several of these needs.

\subsubsection{Engagement, Motivation, and Active Learning}
Traditional training methods can fall short in sustaining high levels of learner engagement and motivation, both of which are pivotal for effective knowledge acquisition and skill development (Woo, 2014). DGBL, by leveraging interactive mechanics, narrative, challenge, and immediate feedback, can significantly boost active participation and foster intrinsic motivation (Tay et al., 2022; Chang et al., 2018). This heightened engagement is frequently proposed as a key factor leading to improved learning outcomes in various contexts, including corporate training (Hainey et al., 2016).

\subsubsection{Safe and Realistic Practice for Complex Skill Development}
Sales success often hinges on the mastery of interpersonal and strategic skills, such as negotiation, objection handling, and adapting communication to diverse client needs (Bussi\`ere, 2017). DGBL can offer simulations of realistic sales scenarios, providing a safe environment in which trainees can practice these complex skills without the real-world consequences of failure (Senderek et al., 2015; Dahalan et al., 2024). This sandboxed approach allows for experimentation, learning from mistakes, and iterative skill refinement, which is crucial for developing robust competencies (Chen et al., 2018).

\subsubsection{Problem-Solving, Critical Thinking, and Learner Agency}
Effective sales professionals are adept problem-solvers who can critically assess situations and devise appropriate strategies. DGBL can be designed to present learners with ill-defined problems or multifaceted challenges that require analytical thinking, strategic decision-making, and creative solutions (Noroozi et al., 2020). Such designs can also enhance learner agency by allowing individuals to make meaningful choices and observe the consequences of those choices within the game environment.

\subsubsection{Scalability, Consistency, and Corporate Upskilling}
Organizations often face substantial training demands as they seek to ensure that sales teams possess current knowledge and skills in response to new products, evolving market conditions, and the general need for continuous professional upskilling (Hern\'andez and Moreno, 2019). DGBL, as a digital modality, offers a scalable and consistent way to deliver high-quality training content across geographically dispersed teams, thereby supporting efficient competency development (Tay et al., 2022).

The strategic integration of DGBL into sales education is therefore proposed not merely as a technological novelty to promote engagement and motivation, but as an informed approach to leverage its unique affordances in addressing organizational training needs, enhancing core sales competencies, and ultimately improving learning and performance outcomes in sales education.

Meta-analytic reviews offer qualified empirical support for these claims, though with important boundary conditions. Wouters et al. (2013) found that serious games outperformed conventional instruction on learning outcomes ($d$ = 0.29) and retention ($d$ = 0.36), with benefits enhanced by multiple practice sessions, collaboration, and supplementary instruction. Merchant et al. (2014), in a meta-analysis of simulation-based learning, found an overall effect size of $d$ = 0.33. However, both reviews focus predominantly on student populations in formal educational settings, and the Merchant analysis in particular draws heavily on classroom and undergraduate contexts. The applicability of these aggregate effects to professional adult learners in sales training---where learner agency, prior domain expertise, and motivation to engage are configured very differently---cannot be assumed without qualification. All et al.'s (2021) systematic review, which gives more attention to workplace and professional DGBL contexts, provides a more cautious aggregate picture and underscores the importance of design quality, instructional alignment, and methodological rigor in determining whether games outperform conventional alternatives. This study situates itself within that more critical tradition.

\subsection{Current Approaches to Sales Education}

\subsubsection{Traditional and Experiential Methods}
Sales education has matured well beyond rudimentary product training and now aims to develop the diverse competencies required in increasingly complex market environments. A variety of pedagogical strategies are employed to bridge theoretical knowledge with practical application, which is crucial for effective knowledge transfer from education to organizational performance (Tho and Trang, 2015).

Traditional methods such as lectures, assigned readings, and classroom discussions continue to play a role in delivering foundational knowledge, including core sales theories, ethical guidelines, and an understanding of market dynamics. However, the dynamic and action-oriented nature of sales warrants a strong focus on experiential learning. As reviewed by Cummins et al. (2013), sales pedagogy widely incorporates techniques intended to offer learners applied experience, such as role-play and simulation. These activities allow learners to practice sales interactions, from initial approaches to closing and objection handling, in a controlled environment, with the aim of building practical skill and often self-efficacy in handling sales situations (Knight et al., 2014).

\subsubsection{Case Study Analysis}
Analyzing real-world or hypothetical business cases helps learners develop diagnostic and strategic thinking skills that are directly relevant to sales challenges.

\subsubsection{Field Projects and Internships}
Direct engagement in sales activities, whether through internships or structured ``selling lab'' projects, provides invaluable hands-on experience and facilitates deeper understanding of the sales process through direct application (Bussi\`ere, 2017).

The emphasis on experiential methods reflects the view that active participation and ``learning by doing'' are critical for skill acquisition and retention in professional disciplines. These approaches are generally intended to foster not only procedural skills, but also critical thinking, problem-solving, and effective interpersonal communication.

Despite the recognized value of these methods, several challenges remain. Transfer from training environments to on-the-job performance can be inconsistent and is shaped by factors in both the training design and the organizational context (Tho and Trang, 2015). Ensuring the quality, consistency, and scalability of resource-intensive experiential methods such as deeply engaging role-play or mentored field experience can also be difficult for educational institutions and corporate training departments. Furthermore, assessing the skills developed through such methods often requires sophisticated rubrics and can remain partly subjective.

Adding to these pedagogical considerations is the ongoing evolution of the sales profession itself. The traditional linear ``seven steps of selling'' model (Moncrief and Marshall, 2005), while still useful, is increasingly supplemented or recontextualized by more dynamic, customer-centric, and digitally mediated approaches, particularly in complex B2B environments (Cardinali et al., 2025). This shift demands that sales education prepare learners for greater complexity, adaptability, and a deeper understanding of value co-creation.

It is in this context that DGBL presents itself as a promising pedagogical innovation. DGBL can encapsulate many principles of experiential learning within interactive and engaging formats. Moreover, it offers potential advantages in terms of scalability, consistency, immediate and individualized feedback, adaptive learning pathways, and the objective capture of rich performance data, thereby addressing some of the limitations of purely traditional methods.

\section{Theoretical and Design Foundations}

\subsection{LM--GM Framework}
In order to deliver the variety of skills and knowledge needed in sales education, a coherent design framework informed by learning theory is necessary in DGBL. Arnab et al. (2015) formulated the learning mechanics--game mechanics framework (LM--GM) by identifying a range of learning mechanics categorized by instructional purpose. The model shows how different game mechanics can be linked to corresponding learning mechanics across increasing levels of cognitive complexity.

\corefigure{lm_gm_framework.pdf}{LM--GM model classified based on Bloom's ordered thinking skills, adapted from Arnab et al. (2015).}{fig:lm-gm-framework}

Correspondingly, the framework outlines numerous game mechanics that can be employed to realize these learning mechanics. Examples include:

\begin{itemize}[leftmargin=1.2cm]
    \item \textbf{Guidance \& Feedback:} Tutorials, hints, pop-up information, and direct feedback messages such as success/failure indicators or explanatory notes.
    \item \textbf{Participation \& Action:} Quests, tasks, levels, puzzles, resource management challenges, and the use of tokens or point systems.
    \item \textbf{Exploration \& Discovery:} Open environments, branching narratives, hidden information elements, and interactive objects.
    \item \textbf{Reflection \& Collaboration:} In-game dialogue systems, forums, team-based challenges, or voting mechanisms.
    \item \textbf{Motivation \& Ownership:} Customizable avatars, leaderboards, collection mechanics, and clear progression systems.
\end{itemize}

The LM--GM framework encourages a design process that begins with clearly defined learning outcomes, followed by the identification of the learning mechanics necessary to achieve them. Appropriate game mechanics are then selected to implement these learning mechanics effectively. This helps ensure that the game's design remains anchored to pedagogical purpose.

For instance, if a learning objective is for sales trainees to classify different types of customer objections, which involves learning mechanics such as identification and discrimination, the game might employ card-sorting challenges or scenario-based quizzes in which players must correctly categorize presented objections to progress. For objectives involving negotiation skills, which rely on experimentation and reflection, the game might use branching dialogue systems with varied outcomes and provide after-action reviews as a feedback mechanic to encourage reflection on choices and consequences.

Furthermore, the LM--GM framework serves not only as a design tool but also as an analytic tool for evaluating existing serious games. It allows researchers and practitioners to deconstruct how specific game features contribute, or fail to contribute, to learning, which is especially valuable when assessing the pedagogical soundness of DGBL interventions. In their own application of the framework to the serious game \textit{Re-Mission}, Arnab et al.\ (2015) note that even within LM--GM-designed games, ``extraneous (pedagogy-independent) game mechanics are often designed to enhance game play,'' with the consequence that ``learning occurs only tangentially.'' This analytic capacity---identifying where game mechanics fail to activate their intended learning functions despite structural alignment---is one reason the framework is adopted as the theoretical lens for the intervention evaluated in this study, and why empirical testing is necessary to verify whether design-level alignment translates into cognitive-level learning gain.

\subsection{Foundations of Game-Based Learning}
Plass, Homer and Kinzer (2015) provide a broader theoretical anchor for game-based learning than design-mapping frameworks alone. They describe game-based learning as resting on four interlocking foundations: cognitive, motivational, affective, and sociocultural. This perspective is useful here because it clarifies that a well-aligned game can still fail if one of those foundations is under-supported. In the present study, LM--GM explains how mechanics were mapped to intended learning functions, while Plass et al. help explain why that design might still produce uneven outcomes across learners.

The present thesis contributes most directly to the \textit{cognitive} foundation in Plass et al.'s model. The central claim is not that game mechanics should simply be aligned with pedagogy, but that their cognitive cost must also be absorbable for the target learner population. The motivational and affective foundations remain visible through the measured roles of engagement and self-efficacy, but they are treated as downstream mechanisms whose expression depends partly on whether learners can first absorb the mechanics inventory without exhausting working-memory capacity. This is where the later Cognitive Load Budget Framework is positioned: not as a rival to Plass's foundations, but as a more operational proposal nested inside the cognitive foundation.

\subsection{Cognitive Load Theory}
As described in the background of this study, cognitive load theory (CLT) is used as one of the main mechanisms that justifies the learning design based on DGBL, in the sense that the mediation of cognitive load through game design is expected to support higher self-efficacy and engagement. CLT begins with the understanding that learning, whatever mechanisms and conditions are used to produce it, depends on the reproduction and organization of knowledge in memory. Yet similar conditions can lead to different outcomes, which raises a cognitive-architectural question: how much working memory is required to learn something at a given time, and how does that relate to the limits of the human cognitive system?

Sweller et al. (2011) describe CLT on the basis that working memory is limited in both capacity and duration, whereas long-term memory is virtually unlimited. The crux of learning lies in transferring information from the temporary space of working memory into the more durable structures of long-term memory. This occurs through schema construction and, over time, through automation.

CLT categorizes cognitive load into three types: intrinsic, extraneous, and germane load. Intrinsic load is the inherent difficulty of the material, influenced by element interactivity. Extraneous load is created by poor instructional design, distracting materials, or unnecessary embellishment and should be minimized as much as possible. Germane load is the mental effort that directly contributes to learning and schema formation. Instruction should therefore aim to maximize this form of effort through strategies that promote reflection, deep processing, and transfer.

These distinctions enable educators and instructional designers to diagnose which type of load learners are experiencing, reduce extraneous load by streamlining content, scaffold intrinsic load to match learner readiness, and amplify germane load to support deeper learning. CLT thus becomes an invaluable tool not only for explaining how learning occurs, but for shaping how it can be facilitated effectively. While useful, CLT also has notable criticisms that should be acknowledged when designing learning materials. First, it depends on a relatively simple representation of memory and may underspecify emotional, episodic, or procedural influences on learning. Second, it risks circularity when failures are retrospectively labelled as extraneous load and successes as germane load. As Popper, discussed in Church (1975), argued more broadly, a theory that can explain everything risks becoming difficult to falsify. Despite these criticisms, CLT remains highly useful in designing instructional media, especially in fields such as sales education, where much knowledge is tacit and only partially formalized.

\subsection{Self-Efficacy}
Digital Game-Based Learning offers a promising avenue for fostering self-efficacy, particularly when it is designed in congruence with principles aimed at managing cognitive load so as to create an optimal environment for skill development and confidence building. As described by Bandura (1977), self-efficacy beliefs are primarily shaped through four sources of information: mastery experiences, vicarious experiences, verbal persuasion, and physiological or affective states.

DGBL interventions can be strategically designed to leverage these sources. For mastery experiences, DGBL can provide repeated opportunities for successful task completion and problem-solving. By segmenting complex sales scenarios into manageable, scaffolded game tasks, learners can accumulate a sequence of scaffolded successes. These achievements build confidence in their own abilities, especially when tasks feel achievable because cognitive load has been managed appropriately.

DGBL can also facilitate vicarious experiences by incorporating models of successful performance, such as non-player characters or simulated peers effectively navigating difficult sales interactions. Similarly, verbal persuasion can be built into game feedback systems through encouraging messages, constructive guidance, and reinforcement for effective play. Finally, DGBL can help shape positive affective states by creating learning environments that are engaging and appropriately challenging without becoming overwhelming. By reducing cognitive overload and learning-related anxiety, and by supporting enjoyment or even flow-like experience (Chang et al., 2018), DGBL can help learners associate the learning process with positive feelings that support stronger efficacy beliefs.

This study therefore posits that a scaffolded DGBL approach, by carefully managing cognitive load and structuring opportunities for success, will provide more favorable conditions for mastery experiences and supportive affective states than traditional digital training methods, which in turn should support stronger self-reported self-efficacy.

\subsection{Engagement}
Alongside cognitive load and self-efficacy, learner engagement stands as a third critical factor influencing the effectiveness of educational interventions, particularly those employing digital and interactive modalities such as DGBL. Engagement in learning is broadly understood as the intensity and quality of a learner's active involvement and commitment to a learning activity or process (Christenson et al., 2012). It is a multifaceted construct, often encompassing behavioral, emotional or affective, and cognitive dimensions (Fredricks et al., 2004), further described below:

\begin{itemize}[leftmargin=1.2cm]
    \item \textbf{Behavioral Engagement:} Observable actions such as participation in tasks, time on task, effort, persistence in the face of difficulty, and adherence to learning protocols.
    \item \textbf{Emotional/Affective Engagement:} Feelings and affective responses to the learning activity, including interest, enjoyment, enthusiasm, and a sense of connection to the task.
    \item \textbf{Cognitive Engagement:} Psychological investment in learning, including focused attention, use of learning strategies, self-regulation, and willingness to grapple with complex ideas and exert mental effort.
\end{itemize}

A substantial body of research indicates a strong positive relationship between learner engagement and desired educational outcomes, including achievement, skill acquisition, and knowledge retention (Hainey et al., 2016; Hern\'andez and Moreno, 2019; Tay et al., 2022). In workplace training, fostering engagement is especially important because adult learners often adopt a pragmatic perspective and seek training that is relevant, stimulating, and clearly applicable to their professional roles (Knowles et al., 2020). Disengagement can lead to superficial learning, poor transfer, and a perception of wasted training investment.

DGBL is frequently cited for its potential to enhance learner engagement compared with more passive instructional methods (Dahalan et al., 2024; Liu and Zhou, 2025). Several characteristics of well-designed games contribute to this potential:

\begin{itemize}[leftmargin=1.2cm]
    \item \textbf{Intrinsic Motivation:} Games often tap into challenge, curiosity, fantasy, and control.
    \item \textbf{Clear Goals and Rules:} Games provide explicit objectives and rules that can improve focus and behavioral engagement.
    \item \textbf{Immediate Feedback:} Continuous feedback helps learners understand the consequences of their decisions and remain cognitively invested (Chang et al., 2020).
    \item \textbf{Sense of Agency and Control:} Games often let players make meaningful choices that affect progression or outcome.
    \item \textbf{Narrative and Storytelling:} Story-driven contexts can enhance emotional and cognitive engagement (Chen et al., 2018).
    \item \textbf{Challenge and Mastery:} Appropriately balanced challenge can motivate persistence and contribute to both affective and behavioral engagement.
\end{itemize}

This study therefore hypothesizes that the proposed DGBL intervention, through game mechanics aligned with learning objectives under the LM--GM framework, will foster higher levels of behavioral, emotional, and cognitive engagement than traditional click-and-continue digital training. This increased engagement is expected to contribute positively to learning outcomes and self-efficacy.

\subsection{Sales Competency}
The ultimate aim of most workplace training initiatives, including the DGBL intervention proposed in this research, is to bring about a positive change in learners' capabilities and subsequently in their job performance. In the domain of sales, this translates to the development and enhancement of sales competency. Competency can be understood as an integrated set of knowledge, skills, abilities, and related personal characteristics that enable effective task performance (Titus, 1994).

Sales competency is multifaceted. While specific emphases vary across roles and industries, several core areas recur in the literature (Cummins et al., 2013): product and market knowledge; sales process skills across prospecting, needs assessment, presentation, objection handling, negotiation, and closing; interpersonal and communication skills; problem-solving and strategic thinking; and self-management attributes such as resilience, proactivity, adaptability, and ethical conduct.

A key challenge in any training program is the valid and reliable assessment of competency development. Traditional assessments often focus on isolated knowledge recall, whereas a more holistic view of competency requires methods that can gauge the integrated application of knowledge and skill in realistic contexts. In DGBL, performance within game scenarios, choices made, strategies employed, and outcomes achieved can all provide meaningful indicators of targeted competency development (All et al., 2021).

The DGBL intervention designed for this study aims to impact several core sales competencies by simulating realistic scenarios and requiring active participation. In the finalized thesis, this broader competency domain is represented most directly through adaptive selling, transfer intention, and related applied sales outcomes reported in the results chapters.

\section{Theoretical Integration}

\subsection{Cognitive Load and Self-Efficacy}
Cognitive Load Theory and Self-Efficacy Theory, while originating from different psychological traditions, describe constructs that are deeply interdependent and reciprocally influential within learning environments, especially those involving complex skill acquisition such as corporate sales training.

The experience of cognitive load during learning can significantly affect self-efficacy. When learners are confronted with tasks that impose excessive extraneous cognitive load because of poorly structured information, confusing interfaces, or an overwhelming density of novel elements, their capacity to process information and engage in meaningful learning is hampered (Sweller et al., 2011). Such situations often lead to frustration, confusion, and a sense of incompetence. These negative learning experiences undermine mastery experiences, which Bandura (1977) identifies as the most powerful source of self-efficacy. Furthermore, excessive cognitive load can induce adverse physiological and affective states, such as anxiety or mental fatigue (Chang et al., 2018), which learners may interpret as evidence of inadequacy.

Conversely, existing self-efficacy can shape the perception and management of cognitive load. Learners with robust self-efficacy tend to approach challenging tasks with greater persistence, effort, and strategic thinking (Bandura, 1986). They may therefore allocate their cognitive resources more effectively and find difficult tasks engaging rather than debilitating. In contrast, learners with low self-efficacy may feel overwhelmed more readily and disengage sooner. Without effective cognitive load management, a downward spiral can occur: high load leads to frustration, which lowers self-efficacy, which in turn reduces persistence and increases the likelihood of further failure.

Instructional designs that actively manage cognitive load are therefore not merely improving information processing efficiency. They are also cultivating conditions more conducive to strong self-efficacy. By making learning tasks feel manageable and achievable, instructional interventions allow learners to experience incremental success and build a portfolio of mastery experiences, thereby supporting confidence and persistence.

The proposed DGBL intervention, with its emphasis on systematic scaffolding and phased increases in task complexity, is intended to interrupt negative cycles of overload and diminished confidence. Instead, it seeks to foster a positive spiral in which manageable cognitive demands support mastery experiences, positive affective states, stronger self-efficacy, greater engagement, and more effective application of learned skills.

\subsection{DGBL as a Cognitive-Cost Tradeoff}
A key premise behind using DGBL in complex professional domains such as sales is that interactivity can organize learning in ways that feel more manageable, meaningful, and confidence-building than a linear presentation. Card-based games in particular can segment decisions, constrain options, and provide immediate contextual feedback, all of which can in principle help learners cope with complexity (Kordaki and Gousiou, 2017). That expectation informed the present design.

At the same time, every mechanic added to a game introduces its own cognitive cost. Learners must decode the interface, understand the token economy, infer the role of each card type, and coordinate those mechanics with the domain content being taught. For some learners this additional inventory may become a productive challenge; for others, especially those with weak prior schemas, it may function as overhead that competes with domain processing. The important design question is therefore not only whether a mechanic is pedagogically aligned, but whether the intended audience can absorb its cost while still learning from it.

This study evaluates that tradeoff under unusually controlled conditions because the game and control arms were deliberately built around the same sales content. Both conditions covered the same sales stages, used the same scenario family, and included analogous quick-check decisions. The main difference was mechanics inventory, not domain scope. This makes the intervention well suited to asking whether a richer mechanics layer improves learning, imposes additional cognitive cost, or does both depending on learner readiness.

\chapter{Conceptual Framework and Hypothesis}
This study is guided by a conceptual framework that posits the scaffolded Digital Game-Based Learning (DGBL) intervention, ``Improve Your Sales Game,'' will influence key learning outcomes by affecting intermediate psychological variables. Specifically, the DGBL intervention, through its theoretically grounded design emphasizing scaffolding, CLT principles, and LM--GM framework application, is expected to lead to lower perceived cognitive load and higher sales self-efficacy and learner engagement compared with a traditional digital training approach. These psychological variables are, in turn, hypothesized to be interrelated and to collectively contribute to the development of adaptive-selling beliefs and a greater intention to apply learning.

\corefigure{conceptual_framework.pdf}{Conceptual framework of the DGBL intervention.}{fig:conceptual-framework}

In the finalized thesis, the broad competency domain described in the original proposal is represented most directly through adaptive selling, transfer intention, and related applied sales outcomes.

Based on this framework and the research questions, the following hypotheses are proposed:

\begin{itemize}[leftmargin=1.2cm]
    \item \textbf{Hypothesis 1.} Participants in the scaffolded DGBL intervention (``Improve Your Sales Game'') will demonstrate more favorable sales training outcomes compared with participants in the traditional click-and-continue digital training group. This hypothesis is supported if the scaffolded DGBL group reports significantly higher learner engagement, higher sales self-efficacy, lower perceived cognitive load, greater improvement in adaptive selling performance, and a higher intention to apply learning.
    \item \textbf{Hypothesis 2.} Within the scaffolded DGBL sales-training condition, perceived cognitive load, learner engagement, and sales self-efficacy will be interrelated and will collectively contribute to adaptive-selling performance in theoretically consistent ways. This hypothesis is supported by significant intercorrelations in the expected directions between perceived cognitive load, sales self-efficacy, and learner engagement, and by positive associations of self-efficacy and engagement with adaptive selling, transfer intention, and related applied sales outcomes.
\end{itemize}

\chapter{Methodology}
\label{sec:method}
This chapter outlines the research methodology employed to address the research questions posed in the introduction of the study. It details the research design, the characteristics of the study participants and sampling strategy, the design of the experimental and control interventions, the instruments used for data collection, the experimental procedure, and the plan for data analysis.

\section{Research Design}
This study utilized a quantitative, experimental research design, specifically a pre-test, post-test randomized controlled trial. This design was selected for its recognized strength in establishing cause-and-effect relationships between an intervention and observed outcomes, and for its ability to minimize threats to internal validity (Kirk, 2013; Shadish et al., 2004).

The RCT involved two parallel groups:

\begin{itemize}[leftmargin=1.2cm]
    \item \textbf{Experimental Group:} Participants in this group received the scaffolded, game-based digital learning sales-training intervention, designed as a card-based digital game.
    \item \textbf{Control Group:} Participants in this group received a traditional click-and-continue digital training module covering the same broad sales concepts, but delivered in a non-gamified, linear format without the interactive scaffolding strategies or game mechanics of the experimental condition.
\end{itemize}

Participants were randomly assigned to either the experimental or control group. Data on key dependent variables, including self-efficacy, perceived cognitive load, learner engagement, transfer intention, and adaptive selling, were collected before the training intervention and immediately after its completion.

\corefigure{rct_design.pdf}{RCT design of DGBL in sales education.}{fig:rct-design}

The choice of an RCT design is justified by its capacity to address the study's primary research question, which seeks to compare the effectiveness of the two training interventions. Random assignment helps ensure baseline equivalence on known and unknown confounding variables, thereby allowing more confident attribution of observed differences to the intervention itself. The study was pre-registered prior to data collection. Behavioral log data confirm that the first participant session was recorded on 31 March 2026 and the final session on 15 April 2026, establishing the data collection window as a two-week period in which all event-level records in the platform database were generated.

\section{Population}
The target population for this study comprised novice sales professionals or individuals aspiring to enter sales roles who had limited prior formal sales training or experience. Because the study was delivered online through a publicly accessible website, participation was open to any adult individual interested in sales training.

Participants were recruited through online channels, including professional networking sites, relevant online forums or communities focused on sales and professional development, and related digital recruitment routes. Eligibility for inclusion in the final analysis primarily depended on completion of the pre-test, the assigned training module, and the post-test measures.

No direct compensation was planned for participation, but participants benefited from free access to a sales training module. Participation was voluntary, and informed consent was obtained digitally before data collection began.

To characterize the sample, demographic information was collected as part of the pre-test questionnaire. This included:

\begin{itemize}[leftmargin=1.2cm]
    \item Age group (e.g., 18--24, 25--34, 35--44, etc.)
    \item Years of experience in sales, if any (e.g., none, less than 1 year, 1--3 years, more than 3 years)
    \item Highest level of education achieved
    \item Familiarity with digital games for learning
    \item Current industry of work, if applicable
\end{itemize}

This demographic data was used descriptively to characterize the sample and to support secondary subgroup interpretation where feasible.

\section{Interventions, Measures, and Procedure}

\subsection{Interventions}
Participants were randomly assigned to one of two digital sales-training interventions: the experimental scaffolded DGBL condition or the control traditional digital training condition. Both interventions covered the same broad sales topics derived from established sales-process models, including prospecting, needs analysis, presentation of solutions, handling objections, and closing. The primary difference lay in instructional approach and delivery format.

\subsubsection{The Scaffolded DGBL Intervention (Experimental Group)}
The experimental intervention was a scaffolded, card-based digital game titled ``Improve Your Sales Game.'' The game was designed on the basis of Cognitive Load Theory to manage cognitive demands and Bandura's Self-Efficacy Theory to foster learner confidence. The core of the game involved a deck of digital cards representing sales scenarios across industries and sales stages. The design process was guided by the LM--GM framework to ensure alignment between learning objectives and game mechanics.

The game unfolded in three distinct phases, each targeting different aspects of the sales process and associated skills:

\textbf{Phase 1: Lead Collection and Retention (Conceptual Understanding \& Identification)}

\begin{itemize}[leftmargin=1.2cm]
    \item \textbf{Learning Objective:} To understand and identify different stages of a sales funnel and classify potential leads.
    \item \textbf{Gameplay:} Participants are presented with event cards describing customer scenarios or lead characteristics and must correctly classify those cards according to predefined sales stages to earn lead tokens. This phase emphasizes foundational knowledge and correct identification, with direct feedback on classification accuracy.
\end{itemize}

\textbf{Phase 2: Sales Process Roleplay (Procedural Skill Application \& Negotiation)}

\begin{itemize}[leftmargin=1.2cm]
    \item \textbf{Learning Objective:} To apply sales techniques in simulated interactions, with a focus on negotiation and objection handling.
    \item \textbf{Gameplay:} Participants use collected lead tokens to engage with simulated customer scenarios. They draw action cards representing sales tactics, communication strategies, or solutions, and must strategically choose plays that successfully navigate the challenges posed by event cards. This phase is intended to provide scaffolded practice, immediate feedback, and mastery experiences.
\end{itemize}

\textbf{Phase 3: Problem-Solving, Evaluation, and Reflection (Strategic Thinking \& Metacognition)}

\begin{itemize}[leftmargin=1.2cm]
    \item \textbf{Learning Objective:} To develop strategic thinking in response to complex or unexpected sales situations and reflect on optimal approaches.
    \item \textbf{Gameplay:} Participants may encounter crisis cards that introduce more complex dilemmas requiring broader strategic judgment. This phase encourages deeper processing, evaluation of alternative strategies, and reflection on the broader sales principles underlying the card-based interactions.
\end{itemize}

The overall design aimed to reduce extraneous cognitive load through clear rules and interface, manage intrinsic load by segmenting the sales process into phases and card-based interactions, and promote germane load through strategic decision-making and reflection, thereby fostering self-efficacy and engagement.

\begin{table}[htbp]
\centering
\caption{A priori LM--GM alignment of the scaffolded DGBL intervention}
\label{tab:lmgm-alignment}
\small
\begin{tabularx}{\textwidth}{L{3.0cm}L{3.0cm}L{2.0cm}Y}
\toprule
\textbf{Game mechanic} & \textbf{Intended learning mechanic} & \textbf{Bloom level} & \textbf{Primary adaptive-sales target} \\
\midrule
Event-card classification & Identification, categorisation, repetition with corrective feedback & Retention / Understanding & Recognising lead cues and matching customer situations to sales stages \\
Immediate accuracy feedback & Reflection, reinforcement, error correction & Understanding & Clarifying why a lead or response fits a given sales stage \\
Lead-token economy & Progression incentive, gated access to later tasks & Retention / Applying & Encouraging accurate foundational decisions before scenario play \\
Phase 2 scenario encounters & Application, decision practice, contextual discrimination & Applying & Selecting responses to objections, needs, and solution-fit situations \\
Action-card choice under constraints & Strategic selection, comparison of alternatives & Applying / Understanding & Choosing tactics appropriate to customer context \\
Crisis / escalation cards & Evaluation, strategic judgment, metacognitive reflection & Understanding / Applying & Responding to difficult or unexpected sales situations \\
\bottomrule
\end{tabularx}
\normalsize
\end{table}

\noindent\textit{Note.} This mapping was specified at the design stage and later used as a retrospective interpretive check on whether the implemented mechanics remained aligned with their intended learning functions in practice.

\subsubsection{The Traditional Digital Training (Control Group)}
Participants in the control group engaged with a digital training module that covered the same core sales topics as the DGBL intervention, but in a non-gamified self-directed format.

\begin{itemize}[leftmargin=1.2cm]
    \item \textbf{Format:} Information was presented through 16 text-based screens with menu-based navigation, allowing learners to move non-linearly, skip familiar material, and revisit earlier sections.
    \item \textbf{Content Delivery:} Core concepts were explained directly through propositional text and static visual support. A Phase 2 scenario component paralleled the game's later encounter structure by presenting branching decision prompts on the same sales topics.
    \item \textbf{Lack of DGBL Elements:} This condition deliberately excluded the card mechanics, token systems, gated game phases, and resource-management structures present in the DGBL intervention.
\end{itemize}

This control condition therefore functioned as a scaffolded self-directed digital module rather than a passive baseline, allowing for a stricter comparison of whether game mechanics added value beyond accessible scenario-based e-learning.

\begin{table}[htbp]
\centering
\caption{Design contrast between the experimental and control conditions}
\label{tab:design-contrast}
\small
\begin{tabularx}{\textwidth}{L{3.1cm}Y Y}
\toprule
\textbf{Design dimension} & \textbf{Scaffolded DGBL condition} & \textbf{Control digital module} \\
\midrule
Content scope & Prospecting, needs analysis, presentation, objections, closing & Same broad adaptive-sales topics \\
Navigation & Mechanic-gated progression across game phases & Menu-based self-navigation across 16 screens \\
Practice format & Card classification, token use, branching encounters & Text-based explanation plus branching scenario prompts \\
Feedback style & Immediate in-game accuracy and outcome feedback & Content review through slides; scenario prompts without resource-management mechanics \\
Learner challenge & Requires decoding rules, tokens, and encounter logic & Requires content selection and scenario judgment without game rules \\
Primary added demand & Mechanics operation plus content processing & Content processing with self-paced navigation \\
\bottomrule
\end{tabularx}
\normalsize
\end{table}

\subsection{Instruments and Measures}
Data for this study were collected using questionnaires administered at pre-test and post-test, alongside behavioral and applied sales indicators. The original measurement design prioritized validated scales where available and theory-informed item development where needed.

\begin{itemize}[leftmargin=1.2cm]
    \item \textbf{Sales Self-Efficacy:} Participants' confidence in performing key sales tasks was assessed with a 10-item scale adapted from Bandura's self-efficacy guidelines and prior sales-education work (Bandura, 1977; Knight et al., 2014).
    \item \textbf{Perceived Cognitive Load:} Subjective cognitive load during the training intervention was assessed post-intervention using a 4-item scale informed by Paas et al. (2008), Chang et al. (2020), and Liu et al. (2023).
    \item \textbf{Learner Engagement:} This multifaceted construct was assessed post-intervention using a 6-item scale developed from Fredricks et al. (2004) and related DGBL research (Woo, 2014).
    \item \textbf{Adaptive Selling Beliefs:} The original proposal specified a scenario-based competency test covering core sales skills; in the finalized thesis, the applied sales domain is represented most directly through adaptive selling beliefs and related post-intervention outcomes.
    \item \textbf{Transfer Intention:} Participants' intention to apply learned sales skills in work contexts was assessed post-intervention using a 3-item scale adapted from Tho and Trang (2015).
\end{itemize}

The study combined these pre-test and post-test self-report instruments with behavioral log data from the training platform. All major scales showed acceptable to excellent internal consistency in the analytic sample.

\begin{table}[htbp]
\centering
\caption{Measurement instruments and internal consistency}
\begin{tabular}{p{6.7cm}ccc}
\toprule
\textbf{Scale} & \textbf{$\alpha$} & \textbf{Items} & \textbf{Rating} \\
\midrule
Pre self-efficacy (1--5) & 0.93 & 10 & Excellent \\
Post self-efficacy (1--5) & 0.96 & 10 & Excellent \\
Post cognitive load (1--7) & 0.77 & 4 & Acceptable \\
Post engagement (1--5) & 0.88 & 6 & Good \\
Post transfer intention (1--5) & 0.85 & 3 & Good \\
Pre ADAPTS-SV (1--7) & 0.91 & 5 & Excellent \\
Post ADAPTS-SV (1--7) & 0.92 & 5 & Excellent \\
\bottomrule
\end{tabular}
\end{table}

\textbf{Self-efficacy} was measured with a 10-item sales self-efficacy scale administered at pre-test and post-test on a 5-point agreement scale, drawing on Bandura's self-efficacy guidelines and prior sales-education work (Bandura, 1977; Knight et al., 2014). \textbf{Adaptive selling} was measured using ADAPTS-SV, a validated 5-item psychometric scale assessing adaptive selling beliefs on a 7-point agreement scale (Spiro and Weitz, 1990; Robinson et al., 2002), administered at both pre-test and post-test. The short-form version was selected deliberately: its brevity minimises additional survey burden and reduces the risk of attrition from survey fatigue, which is a particular concern in voluntary online training contexts. \textbf{Cognitive load} was measured post-test using four items assessing mental effort, perceived difficulty, complexity, and mental strain, following standard subjective load measurement logic (Paas et al., 2003; Paas et al., 2008; Sweller et al., 2011). A validated multidimensional alternative such as Leppink et al.'s (2013) intrinsic/extraneous/germane instrument was not used because the original study instruments had already been finalised for a short online voluntary session, and adding a longer diagnostic battery would have increased survey burden in a design already vulnerable to attrition. \textbf{Engagement} was assessed post-test using six items spanning behavioral, emotional, and cognitive engagement, consistent with multidimensional engagement theory (Fredricks et al., 2004; Woo, 2014). \textbf{Transfer intention} was assessed post-test using three items measuring the intention to apply learned sales skills in future practice, adapted from transfer-focused workplace learning research (Tho and Trang, 2015).

The survey instrument file indicates a 7-point format for cognitive load and a 5-point format for transfer intention, and these codings are consistent with the finalized analysis report. This manuscript therefore follows the finalized report as the authoritative source for scale interpretation and scoring.

\subsection{Procedure}
The study was conducted entirely online. The following steps outline the procedure for participant involvement.

\subsubsection{Recruitment and Initial Contact}
\begin{itemize}[leftmargin=1.2cm]
    \item Participants were recruited through online channels such as professional networking sites, online forums, and related digital communities.
    \item Recruitment materials provided a brief overview of the study's purpose, the voluntary nature of participation, the estimated time commitment, and the benefit of free access to a sales training module.
    \item A link to the study website was provided in the recruitment materials.
\end{itemize}

\subsubsection{Informed Consent and Eligibility}
\begin{itemize}[leftmargin=1.2cm]
    \item Upon accessing the study website, potential participants were presented with an information sheet explaining the research, data usage, confidentiality, withdrawal rights, and researcher contact details.
    \item Participants then provided digital informed consent before proceeding.
    \item A brief screening question was used to ensure that participants were adults.
\end{itemize}

\subsubsection{Pre-Test Data Collection}
\begin{itemize}[leftmargin=1.2cm]
    \item After providing consent, participants completed an online pre-test questionnaire.
    \item The pre-test included demographic information, baseline self-efficacy, and baseline applied sales measures.
\end{itemize}

\subsubsection{Random Assignment}
Upon successful completion of the pre-test, participants were randomly assigned to one of the two intervention conditions:

\begin{itemize}[leftmargin=1.2cm]
    \item Experimental Group: Scaffolded DGBL intervention (``Improve Your Sales Game'').
    \item Control Group: Traditional click-and-continue digital training.
\end{itemize}

The randomization was managed automatically by the online platform hosting the study using a simple assignment algorithm based on time and sequence of access.

\subsubsection{Training Intervention Phase}
Participants were then immediately directed to their assigned digital training module.

\begin{itemize}[leftmargin=1.2cm]
    \item The experimental group played the designed DGBL ``Improve Your Sales Game.''
    \item The control group completed the click-and-continue sales training module.
\end{itemize}

Both interventions were designed to cover the same broad sales content and were intended to be completed in a single session.

\subsubsection{Post-Test Data Collection}
Immediately after completing the assigned intervention, participants completed an online post-test questionnaire. This included post-training self-efficacy, cognitive load, engagement, transfer intention, and adaptive-selling measures.

\subsubsection{Participant Flow and Analytic Sample}
Participant flow is summarized in a CONSORT-style format to clarify recruitment, randomization, attrition, and the final analytic sample.

The platform recorded 179 individuals who accessed the study website. Of these, 124 completed the demographic form and 92 completed the pre-test and were randomized. Allocation favored the game condition numerically because of the observed platform flow during recruitment: 55 participants were assigned to the game condition and 37 to the control condition. Non-completion was substantially higher in the game condition. Twenty-nine of the 55 participants allocated to the game condition did not provide usable post-test data, compared with 8 of the 37 in the control condition. The final per-protocol analytic sample therefore comprised 63 participants: 32 in the game group and 31 in the control group.

This attrition pattern is substantively important. The game group experienced significantly higher attrition than the control group (41.8\% vs. 16.2\%; Fisher's exact test, odds ratio $=0.28$, $p=.010$). In the game arm, dropout was concentrated during phase 2 of the matching game (41\% of game dropouts), phase 1 card classification (24\%), and cases in which participants did not engage meaningfully with the game at all (21\%). In the control arm, most dropouts occurred during slide training (63\% of control dropouts), followed by attrition after the pre-test (25\%).

Baseline self-efficacy did not differ significantly between completers and non-completers, suggesting that attrition was not obviously ability-selective. Nevertheless, the differential attrition implies that the final game sample may represent a more persistent or motivated subset of those initially randomized. Accordingly, the primary analyses reported in this thesis are per-protocol analyses based on completers, with the non-completion pattern interpreted as an outcome-relevant feature of the intervention rather than merely a procedural inconvenience.

\corefigure{consort_flow.png}{CONSORT-style participant flow diagram.}{fig:consort-flow}

\optionalfigure{game_attrition_by_phase.png}{Observed attrition points across the game intervention.}{fig:game-attrition-by-phase}

\section{Data Analysis Plan}
Prior to hypothesis testing, data were screened for accuracy, missing values, and the assumptions of statistical tests. Descriptive statistics were calculated for demographic variables and key outcome measures at pre-test and post-test for both groups. An alpha level of .05 was used for tests of statistical significance.

\begin{table}[htbp]
\centering
\caption{Hypothesis and Analysis Methods}
\begin{tabularx}{\textwidth}{L{3.3cm}L{3.5cm}Y}
\toprule
\textbf{Focus} & \textbf{Variable(s)} & \textbf{Statistical analysis} \\
\midrule
RQ1 / H1 between-group comparisons & Self-efficacy, cognitive load, engagement, transfer intention, and adaptive selling & Between-group comparisons of post-test outcomes and pre-post change scores across conditions. \\
RQ2 / H2 interrelationships within the DGBL group & Perceived cognitive load, learner engagement, sales self-efficacy, and adaptive selling & Pearson correlation coefficients to examine the strength and direction of relationships. \\
Collective contribution of psychological variables & Cognitive load, engagement, self-efficacy predicting adaptive selling & Multiple regression analysis to assess collective and unique contributions. \\
Channeling effect exploration & Cognitive load, self-efficacy, engagement, adaptive selling & Exploratory follow-up interpretation based on observed correlation patterns. \\
\bottomrule
\end{tabularx}
\end{table}

The finalized analysis report treated the final completer sample as the primary analytic set ($N = 63$). For the two primary pre-post outcomes---sales self-efficacy (SE) and adaptive selling (ADAPTS)---between-group comparisons used ANCOVA with the respective pre-test score as a covariate, following CONSORT recommendations for RCTs that include baseline measurement. ANCOVA was preferred over gain-score $t$-tests because it does not impose the implicit constraint that the post-on-pre regression slope equals 1.0, and because partitioning pre-test variance from the error term increases statistical precision even when baseline groups are well-balanced. For the post-test only measures---cognitive load, engagement, and transfer intention---Welch's independent samples $t$-tests were used, as no pre-test baseline was available for these constructs. Within-group pre-post changes on self-efficacy and adaptive selling were tested using paired $t$-tests for descriptive completeness only and are not interpreted as causal evidence in isolation.

Exploratory analyses---including correlation moderation via Fisher $z$ tests, three-cluster readiness profiling with k-means, within-game profile translation using the combined behavioral-and-psychological table format of Chen et al.\ (2018), Tukey HSD pairwise contrasts within the game group, and bootstrapped mediation analysis within the game group---were conducted as theoretically motivated follow-ups to the primary hypotheses. For the within-game cluster follow-up, gain scores (post minus pre) were used for pre-post measures (SE change and ADAPTS gain), post-test scores were used as-is for post-only measures (cognitive load, engagement, transfer intention), and Tukey HSD post-hoc contrasts were computed on the game group only to identify which expertise-cluster pairs drove any omnibus differentiation. These analyses were not pre-registered and are reported as exploratory throughout. Sensitivity checks excluded four participants with extreme completion times (more than 10{,}000 seconds, likely left the browser open); this produced an $n = 59$ sample whose results were consistent in direction and comparable in magnitude to the primary analysis.

\chapter{Results}

\section{Baseline Equivalence}
The two analyzed groups were statistically comparable at baseline. Welch's tests showed no significant differences in pre-test self-efficacy, pre-test adaptive selling, or recorded time on task prior to the outcome analyses. This supports the integrity of randomization among completers, while not eliminating the interpretive importance of differential post-randomization attrition.

\begin{table}[htbp]
\centering
\caption{Baseline equivalence between analyzed groups (0--100 scale)}
\begin{tabular}{p{5.2cm}cccc}
\toprule
\textbf{Measure} & \textbf{Game M (SD)} & \textbf{Control M (SD)} & \textbf{$t$} & \textbf{$p$} \\
\midrule
Pre self-efficacy (\%) & 64.1 (20.0) & 62.6 (17.9) & 0.33 & .746 \\
Pre ADAPTS score (\%) & 63.1 (20.4) & 62.2 (18.4) & 0.20 & .851 \\
Time on task (sec) & 3592 (9624) & 3655 (12333) & -0.02 & .982 \\
\bottomrule
\end{tabular}
\end{table}

\noindent\textit{Note.} Self-efficacy and ADAPTS pre-test scores are reported on a 0--100 scale for consistency with the ANCOVA analyses in \textsection\ref{sec:h1}. Both groups are well-matched at baseline. Because ANCOVA is used regardless of balance, this check serves as descriptive confirmation rather than an analytical gate.

\optionalfigure{baseline_equivalence_plot.png}{Visual comparison of baseline measures between analyzed groups.}{fig:baseline-equivalence-plot}

\section{Hypothesis 1: Between-Group Comparisons}
\label{sec:h1}

Hypothesis 1 predicted that the scaffolded game would produce more favorable post-test outcomes than the conventional digital training format.

\subsection*{Primary Outcomes: ANCOVA (SE and ADAPTS)}

For the two primary pre-post outcomes, ANCOVA with pre-test as covariate provided the primary test. Table~\ref{tab:ancova} shows results with both unadjusted descriptives and ANCOVA-adjusted means.

\begin{table}[htbp]
\centering
\caption{ANCOVA results for primary outcomes (0--100 scale; $n = 63$)}
\label{tab:ancova}
\small
\setlength{\tabcolsep}{4pt}
\resizebox{\textwidth}{!}{%
\begin{tabular}{L{2.6cm}ccccccc}
\toprule
\textbf{Outcome} & \textbf{Game M} & \textbf{Ctrl M} & \textbf{Adj. Game} & \textbf{Adj. Ctrl} & \textbf{$F$(1,60)} & \textbf{$p$} & \textbf{$\eta^2_p$} \\
\midrule
SE Post (\%) & 70.2 (17.5) & 75.0 (13.5) & 69.9 & 75.5 & 3.56 & .064\dag & .056 \\
ADAPTS Post (\%) & 67.9 (19.3) & 72.3 (13.7) & 67.7 & 72.5 & 2.22 & .141 & .036 \\
\bottomrule
\end{tabular}}
\normalsize
\end{table}

\noindent\textit{Note.} Adjusted means estimated at grand mean pre-test score (SE: 63.4; ADAPTS: 62.7). Pre-test covariate $\beta = 0.56$ for both outcomes; model $R^2 = .42$--.46. $\dag p < .10$.

SE showed a marginal trend favouring the control condition ($F$(1,60) = 3.56, $p$ = .064$\dagger$, $\eta^2_p$ = .056), approaching a medium effect size. ADAPTS did not differ significantly between groups ($F$(1,60) = 2.22, $p$ = .141, $\eta^2_p$ = .036, small effect). These results do not support H1 for the full sample, but the SE effect is not trivially small: at $\eta^2_p$ = .056, the study is underpowered to reliably detect or rule out a meaningful group difference. A replication targeting 80\% power at $\alpha$ = .05 would require approximately $n = 120$.

\noindent\textit{ANCOVA interpretation note.} These ANCOVA models provide the aggregate adjusted comparison prespecified for the primary outcomes. However, the continuous moderation analyses reported in Section~\ref{sec:expertise} reveal significant pre-outcome $\times$ condition interactions for both SE and ADAPTS. The ANCOVA estimates should therefore be read as average adjusted effects across heterogeneous slopes rather than as evidence that the treatment effect is constant across all baseline expertise levels.

\subsection*{Sensitivity Analysis: Intention-to-Treat Estimation}

The per-protocol analysis above is restricted to the $N = 63$ participants who completed both the pre-survey and the full training session. To assess robustness to differential attrition (41.8\% game vs.\ 16.2\% control), two intention-to-treat (ITT) sensitivity analyses were conducted.

\textit{ITT-A: no-gain imputation ($n = 80$).} Seventeen non-completers (15 game, 2 control) submitted a pre-survey before discontinuing; their post-test scores were not collected. These participants were included under a conservative no-gain assumption: post-test imputed equal to pre-test (i.e., the training produced zero change for that individual). The remaining 21 non-completers who left before completing the pre-survey were excluded from ITT-A, as no baseline data existed to anchor an imputation.

\textit{ITT-B: worst-case bounding ($n = 101$).} All 38 non-completers were included under extreme assumptions: game non-completers received the lowest observed post-test score on each outcome, control non-completers received the highest. This scenario is an upper-bound exercise --- it represents the most pessimistic possible estimate for the game condition --- rather than a realistic estimate of the true treatment effect.

Table~\ref{tab:itt} reports ANCOVA results for SE and ADAPTS under all three scenarios.

\begin{table}[htbp]
\centering
\caption{ITT sensitivity analysis: ANCOVA under three attrition scenarios}
\label{tab:itt}
\small
\setlength{\tabcolsep}{4pt}
\begin{tabular}{L{4.2cm}C{1.0cm}C{2.0cm}C{1.8cm}C{1.2cm}C{1.2cm}}
\toprule
\textbf{Scenario} & \textbf{$N$} & \textbf{Adj.\ $\Delta$ (G$-$C)} & \textbf{$F$(1,df)} & \textbf{$p$} & \textbf{$\eta^2_p$} \\
\midrule
\multicolumn{6}{l}{\textit{Self-efficacy (post SE $\sim$ condition + pre SE)}} \\
PP: completers only & 63 & $-$5.6\% & $F$(1,61) = 3.56 & .063$\dagger$ & .056 \\
ITT-A: no-gain ($n = 80$) & 80 & $-$5.9\% & $F$(1,78) = 4.69 & .035* & .057 \\
ITT-B: worst-case ($n = 101$) & 101 & $-$20.5\% & $F$(1,99) = 31.94 & $<$.001*** & .244 \\
\midrule
\multicolumn{6}{l}{\textit{Adaptive selling (post ADAPTS $\sim$ condition + pre ADAPTS)}} \\
PP: completers only & 63 & $-$4.9\% & $F$(1,61) = 2.22 & .139 & .036 \\
ITT-A: no-gain ($n = 80$) & 80 & $-$2.8\% & $F$(1,78) = 0.91 & .344 & .012 \\
ITT-B: worst-case ($n = 101$) & 101 & $-$20.6\% & $F$(1,99) = 24.77 & $<$.001*** & .200 \\
\bottomrule
\end{tabular}
\normalsize
\setlength{\tabcolsep}{6pt}
\end{table}

\noindent\textit{Note.} Adj.\ $\Delta$ = adjusted mean difference (game minus control) in percentage-point terms at grand mean pre-test covariate. PP = per-protocol. ITT-A imputes post = pre for 17 partial completers (pre-survey submitted, training not finished). ITT-B places all 38 non-completers at worst/best observed outcomes for game/control respectively. $\dagger p < .10$, $*p < .05$, $***p < .001$.

The ITT-A analysis strengthens the SE finding: the marginal per-protocol trend ($p$ = .063$\dagger$) becomes significant ($p$ = .035*) once the 17 partial completers are included under no-gain imputation. The adjusted mean difference in SE is virtually unchanged ($-$5.6\% vs.\ $-$5.9\%), indicating that the marginal effect in the completer sample was attenuated by a positively selected completer pool rather than inflated by it. The ADAPTS null finding remains non-significant under ITT-A ($p$ = .344), confirming that the per-protocol result is not an artefact of survivor selection. Under ITT-B, both outcomes reach significance, but the extreme effect sizes ($\eta^2_p$ = .24 for SE, .20 for ADAPTS) reflect the mechanical consequence of assigning all no-data non-completers to floor and ceiling values and should not be interpreted as realistic estimates. Taken together, the ITT analyses confirm that the per-protocol findings are conservative: if anything, the true SE disadvantage for the game condition is slightly larger than the completer-only analysis suggests.

\subsection*{Secondary Outcomes: Welch $t$-Tests}

Cognitive load, engagement, and transfer intention (post-test only) were compared using Welch's $t$-tests (Table~\ref{tab:ttests}).

\begin{table}[htbp]
\centering
\caption{Between-group post-test comparisons: secondary outcomes}
\label{tab:ttests}
\small
\begin{tabular}{L{3.4cm}C{2.6cm}C{2.6cm}C{1.1cm}C{1.3cm}C{1.1cm}}
\toprule
\textbf{Outcome} & \textbf{Game M (SD)} & \textbf{Control M (SD)} & \textbf{$t$} & \textbf{$p$} & \textbf{$d$} \\
\midrule
Post cognitive load (1--7) & 4.81 (0.83) & 4.19 (0.93) & 2.83 & .006** & 0.71 \\
Post engagement (1--5) & 3.79 (0.61) & 4.03 (0.61) & -1.54 & .128 & -0.39 \\
Post transfer intention (1--5) & 3.80 (0.69) & 4.01 (0.62) & -1.27 & .209 & -0.32 \\
\bottomrule
\end{tabular}
\normalsize
\end{table}

The only statistically significant between-group difference was cognitive load: the game group reported substantially higher perceived cognitive demand ($t$(59.8) = 2.83, $p$ = .006, $d$ = 0.71) than the control group---a result opposite to the hypothesized direction. Engagement and transfer intention did not differ significantly, though both favoured the control condition with small-to-medium effects.

\subsection*{Interpreting the Cognitive Load Gap}
The higher cognitive load reported by the game group is theoretically important in its own right and does not require a post-hoc item decomposition to be meaningful. The two conditions were intentionally matched on domain content: both covered the same sales stages, used the same scenario family, and targeted the same adaptive-selling concepts. Intrinsic load was therefore designed to be broadly constant across conditions. What differed was the mechanics inventory required to access that content. On that basis, the most plausible interpretation is that the game introduced additional non-intrinsic load relative to the control condition. Whether that additional load was primarily extraneous overhead, partly germane effort, or some mixture of both cannot be cleanly determined from the brief four-item instrument used here.

This measurement limitation should be kept explicit. The cognitive-load scale was selected because the study was conducted in a short voluntary online format where survey burden itself posed an attrition risk. A richer multidimensional instrument such as Leppink et al.'s (2013) intrinsic/extraneous/germane measure would permit stronger subtype claims in future work. For the present study, the robust conclusion is simply that the game required more perceived cognitive effort than the control while not delivering superior aggregate performance, which makes the cognitive-cost side of game design impossible to ignore.

\optionalfigure{between_group_posttest_plot.png}{Between-group comparison of post-test outcomes and ANCOVA-adjusted means.}{fig:between-group-posttest-plot}

\section{Within-Group Change}
Both groups improved significantly from pre-test to post-test on self-efficacy and adaptive selling, confirming that both interventions were educationally active. The control group showed larger absolute gains on both outcomes.

\begin{table}[htbp]
\centering
\caption{Within-group pre-post improvement summary (0--100 scale)}
\small
\begin{tabular}{L{2.1cm}L{2.8cm}L{2.4cm}L{2.4cm}C{1.8cm}}
\toprule
\textbf{Group} & \textbf{Measure} & \textbf{Pre M (SD)} & \textbf{Post M (SD)} & \textbf{Gain} \\
\midrule
Game ($n=32$) & Self-efficacy (\%) & 64.1 (20.0) & 70.2 (17.5) & +6.2** \\
Game ($n=32$) & ADAPTS score (\%) & 63.1 (20.4) & 67.9 (19.3) & +4.8 \\
Control ($n=31$) & Self-efficacy (\%) & 62.6 (17.9) & 75.0 (13.5) & +12.5*** \\
Control ($n=31$) & ADAPTS score (\%) & 62.2 (18.4) & 72.3 (13.7) & +10.1 \\
\bottomrule
\end{tabular}
\normalsize
\end{table}

\noindent\textit{Note.} Within-group SE change: game $p$ = .003**; control $p$ < .001***. Standardized within-group gains: SE game $d$ = 0.52, control $d$ = 0.74; ADAPTS game $d$ = 0.42, control $d$ = 0.54.

\optionalfigure{pre_post_change_plot.png}{Pre- to post-test change in self-efficacy and adaptive selling by condition.}{fig:pre-post-change-plot}

Both interventions produced meaningful pre-post improvement, but the control group's gains were approximately twice as large on self-efficacy. This pattern is consistent with the ANCOVA adjusted means (69.9 vs. 75.5 for SE), where the marginal between-group difference ($p$ = .064$\dagger$) reflects the control group's larger gain rather than a starting-point advantage. The aggregate picture therefore is not that the game harmed learning, but that it was less efficient for this participant mix --- a distinction elaborated by the expertise-conditional analysis in Section~\ref{sec:expertise}.

\section{Hypothesis 2: Interrelationships Within the Game Group}
Hypothesis 2 proposed that cognitive load, engagement, and self-efficacy would relate to one another in theoretically expected ways within the game condition. The results partially supported this hypothesis. The positive relationships among self-efficacy, engagement, transfer intention, and adaptive selling were strong and coherent. However, cognitive load did not show the expected negative relationships with self-efficacy or engagement.

\begin{table}[htbp]
\centering
\caption{Selected correlations within the game group}
\begin{tabularx}{\textwidth}{L{4.0cm}C{1.2cm}C{1.5cm}Y}
\toprule
\textbf{Variable pair} & \textbf{$r$} & \textbf{$p$} & \textbf{Interpretation} \\
\midrule
Cognitive load $\leftrightarrow$ self-efficacy & .043 & .814 & Not significant; opposite to expected direction \\
Cognitive load $\leftrightarrow$ engagement & .236 & .184 & Not significant; opposite to expected direction \\
Self-efficacy $\leftrightarrow$ engagement & .761 & $<.001$ & Strong positive association \\
Self-efficacy $\leftrightarrow$ transfer intention & .918 & $<.001$ & Very strong positive association \\
Engagement $\leftrightarrow$ transfer intention & .833 & $<.001$ & Strong positive association \\
Self-efficacy $\leftrightarrow$ ADAPTS post-test & .877 & $<.001$ & Very strong convergent validity \\
Engagement $\leftrightarrow$ ADAPTS post-test & .736 & $<.001$ & Strong positive association \\
Cognitive load $\leftrightarrow$ ADAPTS post-test & .172 & .340 & Not significant \\
\bottomrule
\end{tabularx}
\end{table}

\optionalfigure{game_group_correlation_heatmap.png}{Correlation matrix for key psychological and outcome variables within the game group.}{fig:game-group-correlation-heatmap}

The strongest message from these correlations is that the motivational and attitudinal constructs formed a tightly connected network inside the game condition. Learners who felt more capable also tended to feel more engaged, more willing to transfer their learning, and more adaptive in their perceived selling approach. What was not supported was the assumption that lower cognitive load would accompany these benefits. This weakens the argument that the game's scaffolding successfully managed mental effort.

\subsection*{Condition Moderation of Construct Relationships (Fisher $z$)}

Fisher $z$ tests compared key bivariate correlations between the game and control groups to assess whether the game restructured the psychological architecture of learning (Table~\ref{tab:fisherz}).

\begin{table}[htbp]
\centering
\caption{Condition moderation of construct relationships (Fisher $z$ tests)}
\label{tab:fisherz}
\small
\begin{tabular}{L{3.8cm}C{1.4cm}C{1.4cm}C{1.6cm}C{1.0cm}C{1.1cm}}
\toprule
\textbf{Relationship} & \textbf{Game $r$} & \textbf{Ctrl $r$} & \textbf{Fisher $z$} & \textbf{$p$} & \textbf{Sig.} \\
\midrule
SE $\times$ ADAPTS & .877*** & .532** & 2.90 & .004 & ** \\
ENG $\times$ ADAPTS & .736*** & .305\dag & 2.37 & .018 & * \\
CL $\times$ ENG & .236 ns & .592*** & -1.66 & .097 & \dag \\
\bottomrule
\end{tabular}
\normalsize
\end{table}

The game condition was associated with significantly tighter coupling between self-efficacy and adaptive selling ($r$ = .877 vs. .532, $p$ = .004): confidence and adaptive selling beliefs were more tightly integrated in game completers than in control participants. Engagement--ADAPTS coupling was similarly amplified ($p$ = .018). Conversely, the strong cognitive load--engagement link present in the control condition ($r$ = .592) was attenuated in the game group ($r$ = .236, $p$ = .097$\dagger$), suggesting that game mechanics may partially decouple engagement from cognitive effort, consistent with self-determination theory: the game can generate intrinsic engagement through its own mechanics rather than only through effortful processing.

\subsection*{Mediation: SE as the Mechanism Linking Engagement to Adaptive Selling}

Bootstrapped mediation analysis ($B$ = 5{,}000) within the game group tested whether SE mediates the effect of engagement on adaptive selling. The indirect effect (ENG $\to$ SE $\to$ ADAPTS) was significant (indirect effect $ab$ = 1.10, 95\% CI [.43, 2.24]). When SE was entered as a mediator, the direct effect of engagement on ADAPTS dropped from $\beta$ = 1.41 to $\beta$ = 0.31 (ns), indicating full mediation. SE alone explained $R^2$ = .77 of adaptive selling variance ($\beta$ = .79, $p$ < .001). No cognitive load mediation pathways were significant. This result should be treated as exploratory given $n$ = 32 and wide confidence intervals, but the direction is theoretically coherent: the game's primary within-group learning mechanism runs through confidence-building rather than cognitive load management.

\section{Behavioral Evidence}
The behavioral logs provide context for interpreting why the game may not have produced superior outcomes. In phase 1, average classification accuracy was 57.9\%, only modestly above a 50\% chance baseline, indicating that the task was genuinely challenging. Performance varied substantially by sales stage, with follow-up and needs analysis performing relatively well but presentation and pre-approach performing poorly. Faster responders tended to be more accurate and reported higher post-test self-efficacy; the correlation between mean phase-1 response time and post-test self-efficacy was $r=-.37$, $p=.031$. This pattern supports the interpretation that slower, more effortful gameplay was linked to lower confidence.

\optionalfigure{phase1_accuracy_by_stage.png}{Phase 1 classification accuracy by sales stage.}{fig:phase1-accuracy-by-stage}

\optionalfigure[0.74\textwidth]{phase1_response_time_vs_self_efficacy.png}{Association between phase-1 response time and post-test self-efficacy in the game group.}{fig:phase1-response-time-self-efficacy}

In phase 2, scenario performance also varied by encounter type. Some event types produced near-perfect win rates, whereas others were much more difficult. Conversion success and win rate correlated positively with engagement, suggesting that behavioral success aligned with a more positive learning experience. However, overall gameplay performance did not consistently predict the main learning outcomes strongly enough to treat in-game success as a reliable proxy for learning.

\optionalfigure{phase2_win_rate_by_encounter.png}{Phase 2 win rate by encounter type.}{fig:phase2-win-rate-by-encounter}

The control module warrants a more careful characterisation than a passive ``click-through'' format suggests. The module was designed with genuine learner agency: participants could navigate content non-linearly via a menu structure, access slides in any order, skip sections they judged to be already understood, and re-visit earlier material. Phase 2 of the control condition consisted of scenario-based decision prompts that directly paralleled the game's Phase 2 encounter structure, applying the same sales content in a branching-choice format. The control was therefore a scaffolded self-directed module with scenario elements, not a purely passive didactic sequence.

Behavioral logs reflect this design. All 31 analyzed control participants viewed all 16 slides; inter-slide timing varied considerably. Twelve participants (39\%) progressed with a mean inter-slide gap below 10 seconds, consistent with skipping material they had prior familiarity with or navigating rapidly before returning, while only six participants (19\%) averaged more than 30 seconds per slide. This heterogeneity in pacing is consistent with agentic navigation rather than uniform disengagement. The control group's equal or better average outcomes in this context therefore reflect a design that offered accessible, self-paced content delivery with genuine interactivity, providing a meaningful comparison to the game rather than a minimal-effort baseline.

\optionalfigure{control_slide_timing_distribution.png}{Distribution of average inter-slide timing in the control condition.}{fig:control-slide-timing-distribution}

\section{Expertise-Conditional Effects}
\label{sec:expertise}

The aggregate between-group comparisons in Section~\ref{sec:h1} obscure substantial heterogeneity across learner expertise levels. To keep the manuscript analytically disciplined, two layers are distinguished. The \textit{predictive} layer asks \textit{for whom} the game is likely to work and is anchored in pre-registered continuous moderation plus baseline-only clustering. The \textit{explanatory} layer then asks \textit{what happened inside the game} for those baseline-defined groups, using a within-game profiling table built from behavioral logs and post-test variables. The predictive layer carries the evidential weight; the explanatory layer makes the pattern easier to interpret.

\subsection*{Continuous Moderation Analysis (Primary)}

OLS regression models of the form \textit{post\_outcome} $\sim$ intercept + condition + pre\_outcome$_\text{centred}$ + condition $\times$ pre\_outcome$_\text{centred}$ were estimated for each primary outcome. A significant interaction term indicates that the game-versus-control difference depends linearly on pre-intervention expertise level --- the defining test for expertise moderation.

Both primary outcomes yielded significant interactions. For adaptive selling performance, the condition $\times$ pre-ADAPTS interaction was $F$(1, 59) = 11.65, $p$ = .0012** ($\hat{\beta}_{\text{int}}$ = 0.541). For self-efficacy, the condition $\times$ pre-SE interaction was $F$(1, 59) = 5.52, $p$ = .022* ($\hat{\beta}_{\text{int}}$ = 0.365). Simple slopes derived from each model characterise the direction and magnitude of the readiness-conditional pattern (Table~\ref{tab:simple-slopes}).

\begin{table}[htbp]
\centering
\caption{Simple slopes: game vs.\ control advantage at three pre-intervention expertise levels}
\label{tab:simple-slopes}
\small
\setlength{\tabcolsep}{4pt}
\begin{tabular}{L{3.6cm}C{2.6cm}C{2.6cm}C{2.6cm}}
\toprule
\textbf{Outcome} & \textbf{Low ($-$1 SD)} & \textbf{Mean (0 SD)} & \textbf{High (+1 SD)} \\
\midrule
ADAPTS (pp) & $-$15.17 & $-$4.83 & +5.51 \\
SE (\%) & $-$0.497 & $-$0.224 & +0.049 \\
\bottomrule
\end{tabular}
\normalsize
\setlength{\tabcolsep}{6pt}
\end{table}

\noindent\textit{Note.} Slope values are the estimated game-minus-control difference in outcome units at each level of pre-intervention expertise. Negative values indicate control advantage; positive values indicate game advantage. ADAPTS is expressed in percentage points; SE in standardised units.

The pattern is monotone for both outcomes: the game disadvantages low-readiness participants, produces near-equivalent outcomes at the mean, and trends positively for high-readiness participants. This is the manuscript's primary answer to the thesis question \textit{for whom does the game work?} It is predictive, decision-relevant, and based on variables available before deployment.

\subsection*{Predictive Baseline Clusters (Supporting)}

An exploratory $k$-means cluster analysis ($k$ = 3, silhouette = 0.49) was conducted on two $z$-standardized \textit{pre-intervention} features only: Pre-SE and Pre-ADAPTS. Restricting clustering to pre-intervention features avoids circularity: including post-test scores or gain variables in the cluster solution would partially define clusters by the very outcomes used to evaluate them. The $k$ = 3 solution was selected via elbow and silhouette criteria; $k$ = 4 isolated a single outlier rather than a theoretically meaningful fourth subgroup. Silhouette = 0.49 indicates acceptable cluster separation on baseline expertise dimensions. Cluster stability was formally verified by re-running $k$-means 100 times with different random seeds: the adjusted Rand index (ARI) across all seed pairs was 1.000 and per-participant label agreement was 100\%, confirming that the solution is perfectly stable and not an artefact of initialisation.

\begin{table}[htbp]
\centering
\caption{Baseline expertise clusters used for readiness profiling ($k = 3$, $n = 63$)}
\begin{tabularx}{\textwidth}{L{2.6cm}C{0.9cm}C{1.6cm}L{2.2cm}L{2.2cm}Y}
\toprule
\textbf{Cluster} & \textbf{$n$} & \textbf{G/C} & \textbf{Pre-SE} & \textbf{Pre-ADAPTS} & \textbf{Interpretation} \\
\midrule
C1 (Novice) & 10 & 4G / 6C & 2.31 (low) & 33\% (low) & Low-readiness entrants \\
C2 (Intermediate) & 28 & 15G / 13C & 3.40 & 57\% & Mid-readiness entrants \\
C3 (Advanced) & 25 & 13G / 12C & 4.17 (high) & 81\% (high) & High-readiness entrants \\
\bottomrule
\end{tabularx}
\end{table}

Within-stratum game vs.\ control comparisons on ADAPTS gain revealed a monotone reversal across expertise levels, consistent with expertise reversal theory (Kalyuga et al., 2003):

\begin{table}[htbp]
\centering
\caption{Within-stratum game vs.\ control: ADAPTS gain by baseline expertise cluster (small-sample corrected effect sizes)}
\label{tab:stratum-effects}
\small
\setlength{\tabcolsep}{4pt}
\begin{tabular}{L{3.8cm}C{1.6cm}C{1.6cm}C{1.4cm}C{2.8cm}C{1.1cm}}
\toprule
\textbf{Expertise stratum} & \textbf{Game gain} & \textbf{Control gain} & \textbf{Hedges' $g$} & \textbf{95\% CI} & \textbf{$p$} \\
\midrule
Low --- C1 (Novice, $n=10$) & +6.7\% & +33.9\% & $-$1.70 & [$-$3.17,\ $-$0.23] & .001 \\
Intermediate --- C2 ($n=28$) & +7.6\% & +9.7\% & $-$0.15 & [$-$0.89,\ +0.60] & n.s. \\
High --- C3 (Advanced, $n=25$) & $-$1.4\% & $-$1.0\% & +0.23 & [$-$0.56,\ +1.02] & n.s. \\
\bottomrule
\end{tabular}
\normalsize
\setlength{\tabcolsep}{6pt}
\end{table}

\noindent\textit{Note.} Hedges' $g$ applies Hedges' (1981) small-sample correction to Cohen's $d$, dividing by $J = 1 - 3/(4df-1)$. 95\% CIs computed via the Hedges--Olkin variance formula. C1 raw $d$ = $-$1.88; correction factor $J$ = 0.903 ($df$ = 8).

Novice learners in the game condition gained substantially less on adaptive selling performance than their control counterparts ($g$ = $-$1.70, 95\% CI [$-$3.17, $-$0.23], $p$ < .001), a large effect consistent with expertise reversal predictions. The wide confidence interval reflects the small C1 stratum ($n$ = 10; $n_\text{game}$ = 4), and the true effect may be larger due to survivor bias in the game novice group. Intermediate learners showed no meaningful group difference ($g$ = $-$0.15, CI spanning zero). Advanced learners showed a small, non-significant edge for the game condition on ADAPTS ($g$ = +0.23, CI spanning zero). To assess construct validity of the cluster solution, the three clusters were cross-tabulated against three demographic variables not used in the clustering. Sales experience showed a significant, monotone association with cluster membership ($\chi^2$(6) = 16.83, $p$ = .001), with mean sales experience increasing from C1 (novice) to C3 (advanced) --- confirming that the pre-intervention expertise clusters map onto an independently measured indicator of relevant prior knowledge. In contrast, game familiarity ($\chi^2$(6) = 1.83, $p$ = .886) and age group ($\chi^2$(6) = 9.41, $p$ = .162) showed no significant association with cluster membership, indicating that the clusters are not confounded by gaming experience or age demographics.

These baseline clusters remain exploratory and should not be treated as confirmatory subgroups. Their role here is supportive: they translate the continuous moderation into an interpretable readiness profile and provide the bridge to the within-game explanatory layer below.

Taken together, the aggregate and readiness-stratified results already suggest that the game and control conditions are not simply stronger or weaker versions of the same instructional pathway. The game imposed higher cognitive cost, and its comparative value depended strongly on learner readiness. The next analyses therefore do not ask again \textit{whether} the game worked on average; they ask \textit{how} the internal pattern of experience differed once learners were inside the game, and whether that helps explain why superficially similar completion of the mechanics led to different downstream value across readiness groups.

\subsection*{Within-Game Profile Translation}

To make the moderation result more tangible, the baseline clusters were then carried into a descriptive analysis of the game arm only. This second layer does \textit{not} answer who should receive the game at entry; rather, it helps explain what happened inside the game once learners from different readiness strata were exposed to the same mechanics. Following Chen et al.\ (2018), behavioral logs and psychological outcomes are displayed together in a single ANOVA table. For readability, the three baseline clusters can be translated descriptively as a low-engagement novice profile (C1), a moderate intermediate profile (C2), and a highly engaged advanced profile (C3). These are explanatory labels derived from post-treatment patterns, not allocation categories.

\begin{table}[htbp]
\centering
\caption{Within-game learning patterns by baseline expertise cluster (game group, $n$ = 32; adapted from Chen et al., 2018, Table 3)}
\label{tab:learning-patterns}
\small
\setlength{\tabcolsep}{4pt}
\begin{tabular}{L{4.2cm}C{2.0cm}C{2.0cm}C{2.0cm}C{1.8cm}}
\toprule
\textbf{Variables} & \textbf{C1: Novice} & \textbf{C2: Intermediate} & \textbf{C3: Advanced} & \textbf{$F$(2,29)} \\
 & \textit{Low-engagement tendency} & \textit{Moderate profile} & \textit{Highly engaged tendency} & \\
 & \textit{$n$ = 4$^\ddagger$} & \textit{$n$ = 15} & \textit{$n$ = 13} & \\
\midrule
\multicolumn{5}{l}{\textit{In-game behavioral performance (from system logs)}} \\[2pt]
Phase 1 accuracy (\%) & 55.00 (20.82) & 54.67 (24.46) & 59.17 (23.53) & 0.128 \\
Phase 2 win rate (\%) & 55.63 (3.27) & 58.07 (12.64) & 62.60 (9.28) & 0.896 \\
Phase 2 conversions & 8.75 (0.50) & 8.73 (1.58) & 9.25 (0.75) & 0.647 \\
Phase 2 failed attempts & 7.00 (0.82) & 6.47 (2.23) & 5.67 (1.67) & 0.960 \\
Phase 2 exact matches & 5.00 (0.82) & 4.80 (2.01) & 5.33 (1.72) & 0.292 \\
Phase 2 redraws & 0.50 (1.00) & 1.47 (3.31) & 0.00 (0.00) & 1.312 \\[4pt]
\multicolumn{5}{l}{\textit{Learning outcomes: gain scores (post $-$ pre)}} \\[2pt]
SE Gain ($\Delta$, 1--5) & 0.78 (0.39) & 0.19 (0.45) & 0.15 (0.44) & 3.311$\dagger$ \\
ADAPTS Gain ($\Delta$\%) & 6.70 (8.60) & 7.57 (15.09) & 1.03 (4.60) & 1.245 \\[4pt]
\multicolumn{5}{l}{\textit{Post-test outcomes (as-is)}} \\[2pt]
Cognitive Load (1--7) & 4.31 (0.97) & 4.80 (0.75) & 4.98 (0.88) & 1.000 \\
Engagement (1--5) & 3.33 (0.24) & 3.65 (0.67) & 4.10 (0.45) & 3.911* \\
Transfer Intention (1--5) & 3.08 (0.42) & 3.56 (0.59) & 4.31 (0.50) & 10.940*** \\
\bottomrule
\end{tabular}
\normalsize
\setlength{\tabcolsep}{6pt}
\end{table}

\noindent\textit{Note.} Cell entries are $M$ ($SD$). $^\ddagger$C1 novice game-group cell $n$ = 4; interpret with caution. Phase 1 accuracy = tokens earned / 10 cards $\times$ 100. Phase 2 win rate = conversions / (conversions + failed attempts) $\times$ 100. Total time excluded: extreme right-skew due to browser-idle outliers (medians: C1 = 32.8 min, C2 = 17.6 min, C3 = 21.8 min; Kruskal--Wallis $H$ = 1.46, $p$ = .483, ns). $\dagger$ $p <$ .10; * $p <$ .05; *** $p <$ .001.

The within-game table reveals a useful dissociation. Objective gameplay metrics were remarkably similar across the three readiness clusters: the groups earned comparable tokens, won comparable proportions of encounters, and differed little on exact matches or failed attempts. What diverged was the psychological meaning of the experience. Engagement and transfer intention increased monotonically from C1 to C3, while self-efficacy gain showed a marginal omnibus trend. In other words, the same mechanics inventory could be traversed on roughly equal behavioral footing while still producing very different subjective and motivational trajectories.

This is where the descriptive labels help. Translating the table into plain language, C1 novices looked like a low-engagement trajectory, C2 like the default moderate trajectory, and C3 like the most engaged trajectory. Those labels are descriptive, not predictive. Their value is interpretive: they show how the predictive readiness model materialised in practice once learners were inside the game. What the game changed most was not raw in-game performance, but the quality of the experience attached to that performance.

To make the explanatory layer more explicit, Tukey HSD pairwise contrasts were then computed on the game group only. These contrasts do not alter the predictive answer about who should receive the game; they refine the explanatory picture by showing which readiness clusters differed once inside the shared game environment.

\begin{table}[htbp]
\centering
\caption{Multiple Comparisons (Tukey HSD) --- Gain-score outcomes within the game group ($n$ = 32)}
\label{tab:tukey-gain}
\footnotesize
\setlength{\tabcolsep}{3.5pt}
\begin{tabular}{L{2.4cm}L{2.4cm}C{2.0cm}C{1.7cm}C{1.2cm}C{1.7cm}C{1.7cm}}
\toprule
\textbf{(I) Cluster} & \textbf{(J) Cluster} & \textbf{Mean Diff (I$-$J)} & \textbf{Std.\ Error} & \textbf{Sig.} & \textbf{95\% CI Lower} & \textbf{95\% CI Upper} \\
\midrule
\multicolumn{7}{l}{\textit{Dependent variable: SE Gain ($\Delta$, 1--5 scale) \quad $F$(2,29) = 3.31, $p$ = .051$\dagger$, MSE = 0.194}} \\[2pt]
C1 Novice
  & C2 Intermediate &$+$0.5883 & 0.1751 & .061$\dagger$ & $-$0.0233 & \hphantom{$-$}1.2000 \\
C1 Novice
  & C3 Advanced & $+$0.6212 & 0.1779 & .050$\dagger$ & $-$0.0003 & \hphantom{$-$}1.2426 \\[2pt]
C2 Intermediate
  & C1 Novice   & $-$0.5883 & 0.1751 & .061$\dagger$ & $-$1.2000 & \hphantom{$-$}0.0233 \\
C2 Intermediate
  & C3 Advanced & $+$0.0328 & 0.1179 & .979 & $-$0.3790 & \hphantom{$-$}0.4447 \\[2pt]
C3 Advanced
  & C1 Novice   & $-$0.6212 & 0.1779 & .050$\dagger$ & $-$1.2426 & \hphantom{$-$}0.0003 \\
C3 Advanced
  & C2 Intermediate &$-$0.0328 & 0.1179 & .979 & $-$0.4447 & \hphantom{$-$}0.3790 \\
\midrule
\multicolumn{7}{l}{\textit{Dependent variable: ADAPTS Gain ($\Delta$\%, 0--100 scale) \quad $F$(2,29) = 1.24, $p$ = .303, MSE = 126.32 [ns]}} \\[2pt]
C1 Novice
  & C2 Intermediate &$-$0.8733 & 4.4722 & .990 & $-$16.4929 & \hphantom{$-$}14.7462 \\
C1 Novice
  & C3 Advanced & $+$5.6692 & 4.5440 & .656 & $-$10.2012 & \hphantom{$-$}21.5397 \\[2pt]
C2 Intermediate
  & C1 Novice   & $+$0.8733 & 4.4722 & .990 & $-$14.7462 & \hphantom{$-$}16.4929 \\
C2 Intermediate
  & C3 Advanced & $+$6.5426 & 3.0115 & .289 & $-$3.9753 & \hphantom{$-$}17.0605 \\[2pt]
C3 Advanced
  & C1 Novice   & $-$5.6692 & 4.5440 & .656 & $-$21.5397 & \hphantom{$-$}10.2012 \\
C3 Advanced
  & C2 Intermediate &$-$6.5426 & 3.0115 & .289 & $-$17.0605 & \hphantom{$-$}3.9753 \\
\bottomrule
\end{tabular}
\normalsize
\setlength{\tabcolsep}{6pt}
\end{table}

\noindent\textit{Note.} Based on observed means. $^*$The mean difference is significant at the .05 level. $\dagger p < .10$. 95\% CIs are simultaneous Tukey confidence intervals controlling familywise error ($\alpha$ = .05; $df_{error}$ = 29, $k$ = 3 groups). C1 novice $n$ = 4; interpret with caution.

\begin{table}[htbp]
\centering
\caption{Multiple Comparisons (Tukey HSD) --- Post-test outcomes within the game group ($n$ = 32)}
\label{tab:tukey-post}
\footnotesize
\setlength{\tabcolsep}{3.5pt}
\begin{tabular}{L{2.4cm}L{2.4cm}C{2.0cm}C{1.7cm}C{1.2cm}C{1.7cm}C{1.7cm}}
\toprule
\textbf{(I) Cluster} & \textbf{(J) Cluster} & \textbf{Mean Diff (I$-$J)} & \textbf{Std.\ Error} & \textbf{Sig.} & \textbf{95\% CI Lower} & \textbf{95\% CI Upper} \\
\midrule
\multicolumn{7}{l}{\textit{Dependent variable: Cognitive Load (1--7) \quad $F$(2,29) = 1.00, $p$ = .380, MSE = 0.686 [ns]}} \\[2pt]
C1 Novice
  & C2 Intermediate &$-$0.4875 & 0.3295 & .554 & $-$1.6381 & \hphantom{$-$}0.6631 \\
C1 Novice
  & C3 Advanced & $-$0.6683 & 0.3347 & .348 & $-$1.8374 & \hphantom{$-$}0.5009 \\[2pt]
C2 Intermediate
  & C1 Novice   & $+$0.4875 & 0.3295 & .554 & $-$0.6631 & \hphantom{$-$}1.6381 \\
C2 Intermediate
  & C3 Advanced & $-$0.1808 & 0.2218 & .834 & $-$0.9556 & \hphantom{$-$}0.5941 \\[2pt]
C3 Advanced
  & C1 Novice   & $+$0.6683 & 0.3347 & .348 & $-$0.5009 & \hphantom{$-$}1.8374 \\
C3 Advanced
  & C2 Intermediate &$+$0.1808 & 0.2218 & .834 & $-$0.5941 & \hphantom{$-$}0.9556 \\
\midrule
\multicolumn{7}{l}{\textit{Dependent variable: Engagement (1--5) \quad $F$(2,29) = 3.91, $p$ = .031*, MSE = 0.308}} \\[2pt]
C1 Novice
  & C2 Intermediate &$-$0.3128 & 0.2207 & .581 & $-$1.0837 & \hphantom{$-$}0.4581 \\
C1 Novice
  & C3 Advanced & $-$0.7683 & 0.2243 & .055$\dagger$ & $-$1.5516 & \hphantom{$-$}0.0150 \\[2pt]
C2 Intermediate
  & C1 Novice   & $+$0.3128 & 0.2207 & .581 & $-$0.4581 & \hphantom{$-$}1.0837 \\
C2 Intermediate
  & C3 Advanced & $-$0.4554 & 0.1486 & .094$\dagger$ & $-$0.9746 & \hphantom{$-$}0.0637 \\[2pt]
C3 Advanced
  & C1 Novice   & $+$0.7683 & 0.2243 & .055$\dagger$ & $-$0.0150 & \hphantom{$-$}1.5516 \\
C3 Advanced
  & C2 Intermediate &$+$0.4554 & 0.1486 & .094$\dagger$ & $-$0.0637 & \hphantom{$-$}0.9746 \\
\midrule
\multicolumn{7}{l}{\textit{Dependent variable: Transfer Intention (1--5) \quad $F$(2,29) = 10.94, $p < .001$, MSE = 0.287}} \\[2pt]
C1 Novice
  & C2 Intermediate &$-$0.4710 & 0.2133 & .278 & $-$1.2160 & \hphantom{$-$}0.2740 \\
C1 Novice
  & C3 Advanced & $-$1.2227$^*$ & 0.2167 & .001 & $-$1.9797 & $-$0.4657 \\[2pt]
C2 Intermediate
  & C1 Novice   & $+$0.4710 & 0.2133 & .278 & $-$0.2740 & \hphantom{$-$}1.2160 \\
C2 Intermediate
  & C3 Advanced & $-$0.7517$^*$ & 0.1436 & .003 & $-$1.2534 & $-$0.2500 \\[2pt]
C3 Advanced
  & C1 Novice   & $+$1.2227$^*$ & 0.2167 & .001 & \hphantom{$-$}0.4657 & \hphantom{$-$}1.9797 \\
C3 Advanced
  & C2 Intermediate &$+$0.7517$^*$ & 0.1436 & .003 & \hphantom{$-$}0.2500 & \hphantom{$-$}1.2534 \\
\bottomrule
\end{tabular}
\normalsize
\setlength{\tabcolsep}{6pt}
\end{table}

\noindent\textit{Note.} Based on observed means. $^*$The mean difference is significant at the .05 level. $\dagger p < .10$. 95\% CIs are simultaneous Tukey confidence intervals controlling familywise error ($\alpha$ = .05; $df_{error}$ = 29, $k$ = 3 groups). Omnibus ANOVA for ADAPTS gain and cognitive load were non-significant; Tukey contrasts are shown for completeness because this layer is explanatory rather than confirmatory. C1 novice $n$ = 4; interpret with caution.

The Tukey HSD contrasts sharpen the explanatory picture. Transfer intention is the most robustly expertise-differentiated within-game outcome: advanced learners (C3) significantly exceeded both novices (C1; $\Delta$ = $+$1.22, $p$ = .001) and intermediate learners (C2; $\Delta$ = $+$0.75, $p$ = .003). Engagement shows the same ordering but with only marginal pairwise contrasts (C3 vs.\ C1: $p$ = .055$\dagger$; C3 vs.\ C2: $p$ = .094$\dagger$), consistent with the smaller omnibus effect and the very small C1 cell. For SE gain, no pairwise contrast crosses the .05 threshold despite the marginal omnibus trend, and neither ADAPTS gain nor cognitive load differentiates the three clusters within the game arm. In short, this explanatory layer is driven less by raw in-game performance than by how motivationally worthwhile the same game felt to different readiness groups.

\optionalfigure{expertise_reversal_heatmap.png}{Effect sizes (Cohen's $d$, game vs.\ control) by outcome and expertise cluster. Red = control advantage; depth of colour reflects magnitude. The monotone reversal pattern on ADAPTS gain across low, intermediate, and high expertise clusters is consistent with expertise reversal theory.}{fig:expertise-reversal-heatmap}

\chapter{Discussion}

\section{Summary of Findings}
This thesis set out to test two linked propositions: first, whether a scaffolded digital game-based intervention would outperform a conventional digital training format in sales training; and second, whether the internal pattern of psychological outcomes inside the game condition would reflect the intended learning logic. Both hypotheses require a qualified answer.

H1 was not supported at the aggregate level. ANCOVA revealed a marginal trend favouring the control group on self-efficacy ($F$(1,60) = 3.56, $p$ = .064$\dagger$, $\eta^2_p$ = .056), with no significant difference on adaptive selling performance. The game produced significantly higher cognitive load than the control ($d$ = 0.71) without corresponding gains in engagement or transfer intention. However, the pre-registered continuous moderation analysis showed that this aggregate comparison averages over a strong readiness gradient: the game was associated with poorer outcomes at low pre-intervention expertise and increasingly competitive outcomes as pre-intervention expertise increased. The exploratory baseline cluster analysis then translated that gradient into concrete learner strata, while the within-game profiling table showed how those strata experienced the same game differently once inside it. The core story is therefore not that the game failed unconditionally, but that its effectiveness appears contingent on learner readiness.

This matters for interpretation because the pattern points to a difference in pathway, not merely a difference in average magnitude. The control condition appears to preserve a more conventional alignment between accessible instruction, confidence gain, and outcome improvement, particularly for novices. The game condition, by contrast, adds cognitive cost and makes readiness more consequential: advanced learners show the strongest transfer-oriented profile, whereas novice learners can traverse broadly similar mechanics without converting that experience into comparable downstream value. The mediation and correlation analyses below are therefore introduced not as isolated follow-ups, but as tests of whether the game reorganized the psychological pathway through which learning value was produced.

H2 received partial support. Within the game group, self-efficacy, engagement, transfer intention, and adaptive selling formed a tightly coherent positive network, corroborated by the finding that self-efficacy fully mediated engagement effects on adaptive selling (indirect effect = 1.10, 95\% CI [.43, 2.24]). Cognitive load did not show the expected negative association with other outcomes, suggesting that subjective load captured a more complicated experience than a simple overload story. Fisher $z$ moderation tests further showed that the game significantly tightened the self-efficacy--ADAPTS coupling ($r$ = .877 vs.\ .532, $p$ = .004), indicating that the game restructured the psychological architecture of learning even when aggregate means were not superior.

The high attrition in the game condition (41.8\% vs.\ 16.2\% in control) is not merely a methodological footnote; it is part of the substantive evaluation. A format that loses a substantial share of participants before completion --- disproportionately around demanding parts of the experience and plausibly over-representing novices --- presents real barriers to deployment, irrespective of completer outcomes.

\section{Interpretation Through Cognitive Load Theory}
The clearest between-group result was higher cognitive load in the game condition, a finding directly opposed to the intervention's original design intent. That result should now be treated as central rather than inconvenient. The two conditions were intentionally matched on content: both covered the same sales stages, used the same scenario family, and targeted the same adaptive-selling domain. Intrinsic load was therefore designed to be broadly constant across conditions. What differed was the mechanics inventory required to access that content. On that basis, the most plausible interpretation is that the game imposed additional non-intrinsic load relative to the control condition.

The present data do not justify a clean empirical decomposition of that added load into extraneous versus germane components. The brief four-item cognitive-load instrument was chosen for reasons of feasibility in a short voluntary online trial, and the earlier post-hoc decomposition attempt does not provide a sufficiently strong basis for subtype claims. The more defensible conclusion is that the game was experienced as cognitively costlier than the control, and that this cost mattered most for learners with the least prior schema to absorb it.

The behavioral data reinforce this reading. Slower processing in the game was associated with lower post-test self-efficacy ($r$ = -.37, $p$ = .031), and phase-1 classification accuracy averaged only 57.9\%, only modestly above chance. Many learners appear to have been still mastering the task structure while they were meant to be engaging with sales content. For novice learners specifically, the game concentrated too much difficulty into the early stages before domain schemas existed to absorb it --- a pattern consistent with expertise reversal predictions (Kalyuga et al., 2003). On theoretical grounds, extraneous overhead is the leading explanation, even if the present instrument cannot isolate it conclusively.

\section{Self-Efficacy, Engagement, and Adaptive Selling Performance}
Within the game group, self-efficacy, engagement, transfer intention, and adaptive selling were all strongly interrelated, indicating that the motivational side of the theoretical framework was functioning coherently. Learners who felt more capable also felt more engaged, more willing to transfer their learning, and stronger in their perceived adaptive selling performance --- consistent with Bandura's (1977, 1986) view that efficacy beliefs are closely tied to persistence, investment, and anticipated future action. In this manuscript, \textit{performance} refers to adaptive selling as operationalized by ADAPTS-SV; the term is therefore shorthand for the study's validated adaptive-selling outcome rather than an observed workplace performance metric.

Critically, mediation analysis showed that self-efficacy \textit{fully} mediated the effect of engagement on adaptive selling (indirect effect = 1.10, 95\% CI [.43, 2.24]), with the direct engagement--ADAPTS path becoming non-significant once self-efficacy was included. The game's primary internal learning mechanism therefore appears to operate through confidence-building: engagement helps only insofar as it first builds self-efficacy, which then drives adaptive selling performance. The Fisher $z$ moderation results extend this interpretation: the game significantly tightened the self-efficacy--ADAPTS coupling relative to the control condition ($r$ = .877 vs.\ .532, $p$ = .004). In the control group, self-efficacy and adaptive selling are correlated but more loosely integrated; in the game group they become much more tightly coupled.

The absence of a meaningful relationship between cognitive load and motivational outcomes (CL $\leftrightarrow$ SE $r$ = .043; CL $\leftrightarrow$ ENG $r$ = .236, both ns) suggests that subjective load captured a complicated experience rather than straightforward overload. Some learners appear to have experienced the game as effortful but worthwhile, whereas others encountered difficulty without a corresponding gain in confidence. The analyzed game sample, after substantial attrition, was also restricted to more persistent participants, which would attenuate observable load--motivation relationships. The between-group result nonetheless stands: the game imposed greater perceived cognitive cost without producing superior aggregate outcomes.

\section{Predictive Readiness and Explanatory Profiles}
The manuscript now distinguishes two legitimate but non-equivalent uses of clustering. The predictive question --- \textit{for whom should the game be deployed?} --- is answered by baseline readiness and continuous moderation. That is the stronger claim because it uses variables available before the intervention. The explanatory question --- \textit{what kind of trajectory did different learners experience once inside the game?} --- is answered by the within-game profiling table, which translates the statistical pattern into a more tangible narrative of low-engagement novice, moderate intermediate, and highly engaged advanced tendencies.

This separation matters methodologically. Post-treatment profile labels are useful for explanation, but they cannot substitute for a predictive allocation rule. The game does not work for learners because they are already ``highly engaged''; rather, higher-readiness learners are the ones most likely to experience the game as engaging and worth transferring. Keeping the predictive and descriptive layers separate avoids circularity while preserving interpretive richness.

The dual-structure analysis also clarifies the role of telemetry. Behavioral logs are not used here as a stand-alone proxy for learning. Instead, they serve as explanatory evidence about how the same mechanics were traversed by learners with different readiness profiles. In that sense, the within-game profiling table is not replacing the predictive model; it is making the predictive model readable.

\section{The Cognitive Load Budget Framework}
The broader theoretical contribution is best understood by placing LM--GM and the present results inside Plass, Homer and Kinzer's (2015) foundations of game-based learning. LM--GM remains the design-mapping framework: it helps specify which mechanics serve which learning functions. Plass et al. provide the wider anchor by distinguishing cognitive, motivational, affective, and sociocultural foundations. The present study contributes most directly to the \textit{cognitive} foundation by proposing a more operational way to think about the cost side of mechanics.

That proposal is framed here as a \textit{Cognitive Load Budget Framework}. In its working form, the framework has three components. First, a game's \textit{raw budget} is the mechanics inventory that the learner must decode in order to access the domain content. Second, prior domain knowledge provides a \textit{schema offset}: the more prior schema a learner brings, the less burdensome that same mechanics inventory becomes. Third, the learner experiences an \textit{effective load} that emerges from the relation between the two. When the effective load remains within the learner's workable range, mechanics can become meaningful challenge; when it exceeds that range, aligned mechanics become overhead.

The present study should be read as a worked example of that logic rather than a definitive validation of it. The game and control conditions were content-matched, yet the game produced a substantial cognitive-load gap ($d$ = 0.71) and a strong readiness interaction. That pattern is consistent with a budget interpretation: the game carried a larger mechanics inventory, and the learners best able to absorb it were those who already possessed stronger prior schema. The framework is therefore defended here at Version 2 --- as an operational and testable mechanism --- rather than at Version 3, which would overclaim direct empirical validation. Its value lies in making explicit something that design frameworks often leave implicit: pedagogical alignment does not guarantee cognitive absorbability.

This reading also clarifies the extension to Arnab et al. (2015). The problem is not that LM--GM fails to align learning and game mechanics. The problem is that aligned mechanics can still generate a \textit{Bloom-level prerequisite dependency}: lower-level mechanics must first be learned before higher-order learning through those mechanics can begin. For advanced learners, that onboarding cost is small. For novices, it may consume the very working memory needed for domain learning. The budget framework therefore adds a readiness-sensitive cognitive screen to the existing alignment logic.

\section{Non-Completers and Deployment}
A central contribution of this manuscript is to keep non-completers visible rather than collapsing the study into the analyzed sample alone. In online intervention studies, attrition is partly a feature of the setting itself: participants can always opt out. That fact should temper over-interpretation. At the same time, attrition patterns are not substantively meaningless. In this study, dropout was higher in the game arm, and the behavioral timing data suggest that at least part of that loss was related to friction within the experience rather than random disappearance alone.

Attrition should therefore be treated in two ways at once. Methodologically, it limits the internal security of a pure completer comparison and requires the ITT sensitivity analyses already reported. Substantively, it is consistent with the broader readiness story: a design that asks too much of low-readiness learners will lose some of them before the putative benefits of the game can activate. This is not proof that every dropout reflects cognitive overload, but it is enough to make non-completion part of the intervention's real-world evaluation.

The practical implication is expertise routing before deployment. The aggregate null on ADAPTS is not a verdict against game-based adaptive sales training; it is a warning against undifferentiated deployment to mixed-readiness audiences in professional-development contexts. For practitioners, the actionable prescription is a pre-training routing step using the same baseline survey data already collected in this study. The ROC-derived threshold remains provisional rather than deployable, because it was derived from a small exploratory stratum, but the logic is plausible: novices need either direct instruction first or a lighter onboarding layer before the full game becomes appropriate.

\section{Limitations}
Several limitations should be acknowledged. First, the final analyzed sample was modest ($N = 63$), and the large post-randomization attrition---especially in the game condition (41.8\% vs.\ 16.2\%)---limits generalizability and means the completer comparison is not equivalent to an intention-to-treat estimate. Post-hoc power analysis (Table~\ref{tab:power}) confirms that the study was adequately powered only for the cognitive load comparison; the two primary outcomes were substantially underpowered. The CL comparison (power = 0.792) is likely to replicate; the SE finding (power = 0.462) and ADAPTS null (power = 0.313) should be treated as preliminary until replicated in larger samples.

\begin{table}[htbp]
\centering
\caption{Retrospective power analysis: required $N$ and achieved power for key outcomes}
\label{tab:power}
\small
\setlength{\tabcolsep}{4pt}
\begin{tabular}{L{3.0cm}C{1.4cm}C{1.4cm}C{1.4cm}C{2.0cm}C{1.8cm}}
\toprule
\textbf{Outcome} & \textbf{Obs.\ $\eta^2_p$ / $d$} & \textbf{$N$} & \textbf{$\alpha$} & \textbf{Power achieved} & \textbf{$N$ for 80\% power} \\
\midrule
Self-efficacy (SE) & .056 & 63 & .05 & 0.462 & $\sim$140 \\
ADAPTS & .036 & 63 & .05 & 0.313 & $\sim$220 \\
Cognitive load & $d=0.71$ & 63 & .05 & 0.792 & $\sim$66 \\
\bottomrule
\end{tabular}
\normalsize
\setlength{\tabcolsep}{6pt}
\end{table}

\noindent\textit{Note.} Power computed using $G^{*}$Power conventions: ANCOVA for SE and ADAPTS (covariate $\rho$ estimated from observed pre-post correlation); independent-samples $t$-test (two-tailed) for CL. $\alpha = .05$ throughout.

Second, the behavioral logs add nuance to the attrition pattern but do not resolve it completely. Without follow-up data from non-completers, the reasons for dropout cannot be confirmed directly. A further interpretive caveat applies to the behavioral telemetry more broadly: as Lim et al. (2021) demonstrate in an online learning context, platform trace data capture only the observable portion of a learner's behavior and do not represent the full range of self-regulatory activity that may occur outside the tracked environment. The first-attempt latency and phase-1 accuracy measures used here should therefore be treated as indicators of in-game processing rather than comprehensive indices of learning engagement. Examining dropout timing adds nuance: participants who exited within 5 minutes of starting (before completing phase 1 content) are more plausibly responding to cognitive friction and task ambiguity than to time unavailability, since 5 minutes is insufficient for a genuine scheduling conflict to resolve. The MAR assumption is strongest for the 5--15 minute dropout band, where participants had engaged with substantive content before discontinuing; these dropouts are most consistent with a time-tolerance-based decision. Participants in the $>$15-minute dropout band had completed a substantial proportion of the task, and their non-completion is more puzzling; if their pre-survey data become available through follow-up, inclusion in a prospective ITT-B analysis would be informative.

Third, the exploratory cluster analysis was not pre-registered. The final cluster solution used two pre-intervention features only (Pre-SE and Pre-ADAPTS), yielding silhouette = 0.49 --- acceptable but not strong separation. This pre-intervention-only approach was adopted to avoid circularity in cluster derivation: an earlier exploratory version of the analysis used 10 features including post-test and gain variables (silhouette = 0.26), which risks partially defining clusters by the very outcomes used to evaluate them; that version was replaced in the final analysis. Stability of the final $k$ = 3 solution \textit{was} formally tested across 100 random seeds and yielded ARI = 1.000, indicating that the solution was not an artefact of initialisation. That said, the expertise-reversal finding ($d$ = -1.88 for novices) remains exploratory and should be confirmed in a pre-registered replication with a fully powered sample before being treated as definitive. A survivor-bias caveat also applies: game-arm novices who completed the study ($n$ = 4) may be the most persistent subset of novice participants, meaning the true game disadvantage for novices could be larger than observed.

Fourth, all primary outcome measures are self-report. The high SE--ADAPTS correlation ($r$ = .877 in the game group) is partly attributable to construct overlap---both measures tap confidence in sales tasks---and shared method variance from the same post-test session. No objective performance measure was included; all ADAPTS-related findings should be interpreted as self-report beliefs about adaptive selling, not demonstrated skill acquisition.

Fifth, the study focused on short-term post-test outcomes and perceived transfer intention, without long-term follow-up on retention or real-world skill transfer. The transfer intention measure (Tho and Trang, 2015) captures post-training intention rather than demonstrated on-the-job transfer; intention-behavior gaps are well-documented in the training transfer literature (Blume et al., 2010), and transfer intention should be interpreted accordingly as a proximal motivational indicator rather than an outcome measure.

Sixth, a format--measurement alignment concern warrants acknowledgement as an alternative explanation for the ADAPTS null finding. The control condition delivered content in propositional text slides, and the ADAPTS-SV questionnaire assesses adaptive selling beliefs through propositional, declarative statements about sales behavior (e.g., ``I vary my sales style from situation to situation''). The surface format of the measurement instrument therefore more closely matches the surface format of the control intervention than the game format. If any portion of ADAPTS-SV performance reflects familiarity with propositional descriptions of selling concepts, independent of genuine belief change, the control group's advantage may be partially attributable to this format match rather than to superior learning. This alternative explanation cannot be ruled out from the present data. Future studies comparing game-based and text-based training should consider including scenario-based or process-oriented performance measures that are format-neutral relative to both intervention arms.

Seventh, a measurement timing concern applies to comparative GBL research more broadly, and to this study in particular. Immediate post-test comparisons are structurally favourable to instructional conditions that encode propositions directly: a text-based module that presents propositional content is compared, moments later, on a measure of propositional knowledge or propositional beliefs, producing a near-identical format--timing match. Game-based learning, by contrast, operates on different psychological levers --- schema formation through repeated simulation, self-efficacy building through mastery experiences, and motivational priming for continued post-training engagement. These mechanisms are not primarily propositional, and their outputs are not primarily immediate: they are likely to manifest more clearly in delayed retention, in on-the-job behavioural transfer, and in self-directed post-training learning than in immediate post-test measurement. This study attempted to partially address the problem by using psychometric belief measures (SE, ADAPTS) rather than declarative knowledge tests, avoiding the most direct format match. However, even belief measures at immediate post-test capture the propositional-encoding advantage of the control condition more directly than the game's motivational--schema pathway. Empirical support for this temporal asymmetry comes from Wouters et al.'s (2013) meta-analysis, which found that the serious game advantage was larger for \textit{retention} ($d$ = 0.36) than for \textit{immediate learning} ($d$ = 0.29) --- a pattern consistent with games producing learning through mechanisms that compound over time rather than being fully expressed at the moment of training completion. The SE-mediated engagement pathway observed within our game group is consistent with precisely this mechanism: the game built confidence and motivational priming that, measured at a delayed follow-up, would be predicted to generate continued self-directed learning and behavioural change. We captured the mechanism but not its product. Future studies in this space should treat delayed measurement --- at a minimum 4--6 weeks post-training, with an on-the-job transfer indicator --- as a primary rather than supplementary endpoint, and should frame immediate post-test comparisons accordingly as conservative lower-bound estimates of game effectiveness rather than definitive outcome assessments.

Eighth, the sample composition limits generalizability. The study recruited participants from a single professional development context, with approximately 69\% working in the education sector and the remainder distributed across industries, predominantly from Southeast Asia, in a voluntary online enrolment setting. This sample profile is unlikely to be representative of B2B sales professionals in other sectors, regions, or organizational training contexts. The expertise distribution (with a relatively small novice stratum, $n$ = 10) may also differ from the target population in deployed training contexts, where novices are often the primary intended audience.

Finally, the study was conducted in a web-based self-enrolment environment, which improves ecological realism for digital training but reduces control over attention and context. Regarding the comparison group: the control module included genuine learner agency --- non-linear menu navigation, content skipping for familiar material, and a Phase 2 scenario component structurally parallel to the game's Phase 2. The heterogeneous pacing observed in behavioral logs (39\% of control participants with sub-10-second inter-slide gaps) is consistent with agentic selective review rather than disengagement, and should not be interpreted as evidence that the control was a minimal-effort baseline. If anything, this design makes the comparison more ecologically valid; future studies might nonetheless add explicit measures of within-module navigation behaviour to characterise control condition engagement more precisely.

\section{Future Research}
The most important next step is a pre-registered replication that stratifies participants by prior sales knowledge before randomisation, enabling a confirmatory test of the expertise-conditional hypothesis. A sufficiently powered study (approximately $n = 120$ for 80\% power at $\eta^2_p$ = .056) with formal expertise moderator analysis---replacing the post-hoc cluster approach with a pre-registered continuous moderation model (pre-ADAPTS $\times$ condition interaction)---would provide stronger evidence about for whom game-based sales training works and under what conditions.

Future designs should also test redesigned game versions that separate interface onboarding from instructional content, with a dedicated non-scored tutorial phase before domain-relevant gameplay begins. Dynamic difficulty adjustment mechanisms---triggered by phase-1 accuracy or response time thresholds---could be tested as a way to route novice learners toward more supported instruction without abandoning the game format entirely (Serrano-Laguna et al., 2012; Zohaib, 2018).

Longer-term follow-up is needed to determine whether the expertise-conditional pattern persists, attenuates, or reverses over time. More fundamentally, delayed measurement should be treated as a primary endpoint in future GBL comparative trials. Immediate post-test comparisons are structurally conservative for game conditions: games operate through schema-building, confidence priming, and motivational carry-forward mechanisms whose outputs compound post-training rather than fully expressing at the moment of completion. A 4--6 week follow-up with behavioural transfer indicators and retention measures would provide a more ecologically valid test of whether the game's motivational architecture translates into durable learning advantage, and would directly test the delayed-effect prediction that our SE-mediation finding implies. Finally, adopting validated multidimensional cognitive load instruments (e.g., Leppink et al.'s 2013 intrinsic/extraneous/germane subscale) in future studies would allow a direct test of which load components drive expertise-conditional effects, instead of inferring them indirectly from a brief global measure.

\chapter{Conclusion}
This thesis evaluated a scaffolded digital game-based intervention for adaptive sales training in a professional-development context against a conventional digital module in an online randomised controlled trial. At the aggregate level, the game did not produce superior outcomes on self-efficacy or adaptive selling performance, and it imposed substantially higher cognitive load than the control ($d$ = 0.71). Read at that level alone, the trial is a null-to-negative result. Read through the pre-registered moderation analysis, however, a more useful pattern emerges: the game was associated with worse outcomes for low-readiness learners and increasingly competitive outcomes as readiness increased.

That predictive readiness pattern is the thesis's main answer to the question \textit{for whom does game-based sales training work?} The answer is not ``for the highly engaged'' in any predictive sense, because engagement is visible only after learners have entered the game. Rather, the game works most plausibly for learners who arrive with enough prior schema to absorb the mechanics inventory without exhausting the working memory needed for domain learning. The baseline clusters and the within-game profiling table then make that result readable: novices tended to experience the same mechanics as a lower-value trajectory, while advanced learners were more likely to experience them as engaging and transferable.

The clearest empirical signal enabling this interpretation is the cognitive-load gap. Because the game and control conditions were intentionally matched on sales content, the higher perceived load in the game condition is most plausibly attributable to the additional mechanics layer rather than to higher intrinsic content complexity. The present instrument does not allow a clean empirical separation of extraneous and germane load, so that subtype claim remains open. What the data do show robustly is that the cognitive-cost side of game design matters, and that it interacts with learner readiness rather than averaging out harmlessly.

This is where the study's theoretical contribution sits. Arnab et al.'s LM--GM framework helps designers ask whether a game mechanic is aligned to a learning function. Plass, Homer and Kinzer's (2015) foundations of game-based learning remind us that alignment is only one part of the story; games also rest on cognitive, motivational, affective, and sociocultural foundations. The present thesis extends the \textit{cognitive} foundation by proposing a Cognitive Load Budget Framework: a design heuristic in which mechanics inventory creates a raw budget, prior schema supplies a learner-specific offset, and effective load determines whether aligned mechanics function as productive challenge or as overhead. The current study should be treated as a worked example of that logic, not as its final validation. Its value lies in turning a familiar but often implicit design intuition into a testable proposition.

The motivational findings complement rather than contradict that account. Within the game group, self-efficacy fully mediated the effect of engagement on adaptive selling performance, and the game significantly tightened the self-efficacy--ADAPTS relationship relative to the control condition. In other words, the game did contain a viable internal learning mechanism, but that mechanism appears to have activated mainly among learners who persisted long enough and entered with enough readiness to benefit from it.

Practically, the study supports a more selective deployment logic for digital game-based sales training. The lesson is not that game-based training should be abandoned, but that it should not be deployed indiscriminately to mixed-readiness audiences. Expertise routing, lighter onboarding for novices, and dynamic difficulty adjustment become design necessities rather than optional polish when mechanics cost is high. The attrition pattern strengthens that conclusion: even in a voluntary online context where some dropout is unavoidable, a format that disproportionately loses lower-readiness learners before completion must count that loss as part of its intervention effect.

The study therefore ends with a narrower but stronger claim than the one it began with. Scaffolded game-based sales training is not a universal upgrade over conventional e-learning. It is a precision intervention whose value depends on whether the audience can absorb the mechanics inventory that carries the pedagogy. That conditionality is not a weakness of the thesis. It is its central contribution.

\begin{thebibliography}{99}

\bibitem{all2021}
All, A., Castellar, E. N. P., and Van Looy, J. (2021).
Digital game-based learning effectiveness assessment: Reflections on study design.
\textit{Computers \& Education, 167}, 104160.

\bibitem{arnab2015}
Arnab, S., Lim, T., Carvalho, M. B., Bellotti, F., de Freitas, S., Louchart, S., Suttie, N., Berta, R., and De Gloria, A. (2015).
Mapping learning and game mechanics for serious games analysis.
\textit{British Journal of Educational Technology, 46}(2), 391--411.

\bibitem{plass2015}
Plass, J. L., Homer, B. D., and Kinzer, C. K. (2015).
Foundations of game-based learning.
\textit{Educational Psychologist, 50}(4), 258--283.

\bibitem{blume2010}
Blume, B. D., Ford, J. K., Baldwin, T. T., and Huang, J. L. (2010).
Transfer of training: A meta-analytic review.
\textit{Journal of Management, 36}(4), 1065--1105.

\bibitem{bandura1977}
Bandura, A. (1977).
Self-efficacy: Toward a unifying theory of behavioral change.
\textit{Psychological Review, 84}(2), 191--215.

\bibitem{bandura1986}
Bandura, A. (1986).
\textit{Social Foundations of Thought and Action: A Social Cognitive Theory}.
Prentice Hall.

\bibitem{barling1983}
Barling, J., and Beattie, R. (1983).
Self-efficacy beliefs and sales performance.
\textit{Journal of Organizational Behavior Management, 5}(1), 41--51.

\bibitem{bussiere2017}
Bussi\`ere, D. (2017).
Understanding the sales process by selling.
\textit{Marketing Education Review, 27}(2), 86--91.

\bibitem{cardinali2025}
Cardinali, S., Giovannetti, M., Hautam\"aki, P., and Burgdorff Jensen, K. (2025).
A competency-based curriculum content design for complex B2B sales in higher education institutions.
\textit{Journal of Marketing Education}. Advance online publication.

\bibitem{chang2018}
Chang, C.-C., Warden, C. A., Liang, C., and Lin, G.-Y. (2018).
Effects of digital game-based learning on achievement, flow and overall cognitive load.
\textit{Australasian Journal of Educational Technology, 34}(4).

\bibitem{chang2020}
Chang, C.-Y., Kao, C.-H., Hwang, G.-J., and Lin, F.-H. (2020).
From experiencing to critical thinking: A contextual game-based learning approach to improving nursing students' performance in electrocardiogram training.
\textit{Educational Technology Research and Development, 68}(3), 1225--1245.

\bibitem{cummins2013}
Cummins, S., Peltier, J. W., Erffmeyer, R., and Whalen, J. (2013).
A critical review of the literature for sales educators.
\textit{Journal of Marketing Education, 35}(1), 68--78.

\bibitem{dahalan2024}
Dahalan, F., Alias, N., and Shaharom, M. S. N. (2024).
Gamification and game-based learning for vocational education and training: A systematic literature review.
\textit{Education and Information Technologies, 29}(2), 1279--1317.

\bibitem{evans2024}
Evans, P., Vansteenkiste, M., Parker, P., Kingsford-Smith, A., and Zhou, S. (2024).
Cognitive load theory and its relationships with motivation: A self-determination theory perspective.
\textit{Educational Psychology Review, 36}(1), 7.

\bibitem{fredricks2004}
Fredricks, J. A., Blumenfeld, P. C., and Paris, A. H. (2004).
School engagement: Potential of the concept, state of the evidence.
\textit{Review of Educational Research, 74}(1), 59--109.

\bibitem{hainey2016}
Hainey, T., Connolly, T. M., Boyle, E. A., Wilson, A., and Razak, A. (2016).
A systematic literature review of games-based learning empirical evidence in primary education.
\textit{Computers \& Education, 102}, 202--223.

\bibitem{hernandez2019}
Hern\'andez, M., and Moreno, J. (2019).
A systematic literature review on organizational training using game-based learning.
In V. Agredo-Delgado and P. H. Ruiz (Eds.), \textit{Human-Computer Interaction} (pp. 1--18).
Springer.

\bibitem{hung2014}
Hung, C.-M., Huang, I., and Hwang, G.-J. (2014).
Effects of digital game-based learning on students' self-efficacy, motivation, anxiety, and achievements in learning mathematics.
\textit{Journal of Computers in Education, 1}(2), 151--166.

\bibitem{kalyuga2003}
Kalyuga, S., Ayres, P., Chandler, P., and Sweller, J. (2003).
The expertise reversal effect.
\textit{Educational Psychologist, 38}(1), 23--31.

\bibitem{kirschner2006}
Kirschner, P. A., Sweller, J., and Clark, R. E. (2006).
Why minimal guidance during instruction does not work: An analysis of the failure of constructivist, discovery, problem-based, experiential, and inquiry-based teaching.
\textit{Educational Psychologist, 41}(2), 75--86.

\bibitem{knight2014}
Knight, P., Mich, C. C., and Manion, M. T. (2014).
The role of self-efficacy in sales education.
\textit{Journal of Marketing Education, 36}(2), 156--168.

\bibitem{kordaki2017}
Kordaki, M., and Gousiou, A. (2017).
Digital card games in education: A ten-year systematic review.
\textit{Computers \& Education, 109}, 122--161.

\bibitem{paas1994}
Paas, F. G. W. C., and van Merri\"{e}nboer, J. J. G. (1994).
Instructional control of cognitive load in the training of complex cognitive tasks.
\textit{Educational Psychology Review, 6}(4), 351--371.

\bibitem{paas2003}
Paas, F., Tuovinen, J. E., Tabbers, H., and Van Gerven, P. W. M. (2003).
Cognitive load measurement as a means to advance cognitive load theory.
\textit{Educational Psychologist, 38}(1), 63--71.

\bibitem{mayer2009}
Mayer, R. E. (2009).
\textit{Multimedia Learning} (2nd ed.).
Cambridge University Press.

\bibitem{moncrief2005}
Moncrief, W. C., and Marshall, G. W. (2005).
The evolution of the seven steps of selling.
\textit{Industrial Marketing Management, 34}(1), 13--22.

\bibitem{noroozi2020}
Noroozi, O., Dehghanzadeh, H., and Talaee, E. (2020).
A systematic review on the impacts of game-based learning on argumentation skills.
\textit{Entertainment Computing, 35}, 100369.

\bibitem{reeve2012}
Reeve, J. (2012).
A self-determination theory perspective on student engagement.
In S. L. Christenson, A. L. Reschly, and C. Wylie (Eds.), \textit{Handbook of Research on Student Engagement} (pp. 149--172).
Springer.

\bibitem{senderek2015}
Senderek, R., Brenken, B., and Stich, V. (2015).
The implementation of game-based learning as part of corporate competence development.
In \textit{2015 International Conference on Interactive Collaborative and Blended Learning} (pp. 44--51).

\bibitem{sweller2011}
Sweller, J., Ayres, P., and Kalyuga, S. (2011).
\textit{Cognitive Load Theory}.
Springer.

\bibitem{tay2022}
Tay, J., Goh, Y. M., Safiena, S., and Bound, H. (2022).
Designing digital game-based learning for professional upskilling: A systematic literature review.
\textit{Computers \& Education, 184}, 104518.

\bibitem{tho2015}
Tho, N. D., and Trang, N. T. M. (2015).
Can knowledge be transferred from business schools to business organizations through in-service training students? SEM and fsQCA findings.
\textit{Journal of Business Research, 68}(6), 1332--1340.

\bibitem{woo2014}
Woo, J.-C. (2014).
Digital game-based learning supports student motivation, cognitive success, and performance outcomes.
\textit{Journal of Educational Technology \& Society, 17}(3), 291--307.


\bibitem{bosche2011}
B\"osche, W., and Kattner, F. (2011).
Fear of (serious) digital games and game-based learning?: Causes, consequences and a possible countermeasure.
\textit{International Journal of Game-Based Learning, 1}(3), 1--15.

\bibitem{chen2018}
Chen, Z.-H., Chen, H. H.-J., and Dai, W.-J. (2018).
Using narrative-based contextual games to enhance language learning.
\textit{Journal of Educational Technology \& Society, 21}(3), 186--198.

\bibitem{christenson2012}
Christenson, S. L., Reschly, A. L., and Wylie, C. (Eds.). (2012).
\textit{Handbook of Research on Student Engagement}.
Springer.

\bibitem{church1975}
Church, A. (1975).
Karl R. Popper: The logic of scientific discovery.
\textit{Journal of Symbolic Logic, 40}(3), 471--471.

\bibitem{kirk2013}
Kirk, R. E. (2013).
\textit{Experimental Design: Procedures for the Behavioral Sciences} (4th ed.).
Sage.

\bibitem{knowles2020}
Knowles, M. S., Holton III, E. F., Swanson, R. A., and Robinson, P. A. (2020).
\textit{The Adult Learner: The Definitive Classic in Adult Education and Human Resource Development} (9th ed.).
Routledge.

\bibitem{liu2023}
Liu, Y.-C., Wang, W.-T., and Huang, W.-H. (2023).
The effects of game quality and cognitive loads on students' learning performance in mobile game-based learning contexts: The case of system analysis education.
\textit{Education and Information Technologies, 28}(12), 16285--16310.

\bibitem{leppink2013}
Leppink, J., Paas, F., Van der Vleuten, C. P. M., Van Gog, T., and Van Merri\"enboer, J. J. G. (2013).
Development of an instrument for measuring different types of cognitive load.
\textit{Behavior Research Methods, 45}(4), 1058--1072.

\bibitem{liu2025}
Liu, S., and Zhou, H. (2025).
Does gamified training improve training performance? A dual-pathway moderated mediation model.
\textit{Journal of Business Research, 188}, 115086.

\bibitem{deci2000}
Deci, E. L., and Ryan, R. M. (2000).
The ``what'' and ``why'' of goal pursuits: Human needs and the self-determination of behavior.
\textit{Psychological Inquiry, 11}(4), 227--268.

\bibitem{paas2008}
Paas, F., Ayres, P., and Pachman, M. (2008).
Assessment of cognitive load in multimedia learning: Theory, methods and applications.
In J. L. Plass, R. Moreno, and R. Br\"unken (Eds.), \textit{Cognitive Load Theory} (pp. 11--35).
Cambridge University Press.

\bibitem{shadish2004}
Shadish, W., Cook, T., and Campbell, D. (2004).
\textit{Quasi-Experimental Designs for Generalized Causal Inference}.
Houghton Mifflin.

\bibitem{sousa2014}
Sousa, M., and Costa, E. (2014).
Game-based learning improving leadership skills.
\textit{EAI Endorsed Transactions on Serious Games, 1}(3), e2.

\bibitem{titus1994}
Titus, A. A. (1994).
Competence at Work, by Lyle M. Spencer, Jr., and Signe M. Spencer.
\textit{Human Resource Development Quarterly, 5}(4), 391--395.

\bibitem{lim2021}
Lim, L. A., Dawson, S., Gašević, D., Joksimović, S., Pardo, A., Fudge, A., and Gentili, S. (2021).
Impact of learning analytics feedback on self-regulated learning: Triangulating behavioural logs with students' recall.
In \textit{Proceedings of the 11th International Learning Analytics and Knowledge Conference} (LAK21) (pp. 364--374).
ACM.

\bibitem{wouters2013}
Wouters, P., Van Nimwegen, C., Van Oostendorp, H., and Van der Spek, E. D. (2013).
A meta-analysis of the cognitive and motivational effects of serious games.
\textit{Journal of Educational Psychology, 105}(2), 249--265.

\bibitem{merchant2014}
Merchant, Z., Goetz, E. T., Cifuentes, L., Keeney-Kennicutt, W., and Davis, T. J. (2014).
Effectiveness of virtual reality-based instruction on students' learning outcomes in K-12 and higher education: A meta-analysis.
\textit{Computers \& Education, 70}, 29--40.

\bibitem{merrill2002}
Merrill, M. D. (2002).
First principles of instruction.
\textit{Educational Technology Research and Development, 50}(3), 43--59.

\bibitem{gagne1985}
Gagn\'{e}, R. M. (1985).
\textit{The Conditions of Learning} (4th ed.).
Holt, Rinehart and Winston.

\bibitem{robinson2002}
Robinson, L., Marshall, G. W., Moncrief, W. C., and Lassk, F. G. (2002).
Toward a shortened measure of adaptive selling.
\textit{Journal of Personal Selling \& Sales Management, 22}(2), 111--118.

\bibitem{serrano2012}
Serrano-Laguna, A., Torrente, J., Moreno-Ger, P., and Fernández-Manjón, B. (2012).
Tracing a little for big improvements: Application of learning analytics and videogames for student assessment.
\textit{Procedia Computer Science, 15}, 203--209.

\bibitem{spiro1990}
Spiro, R. L., and Weitz, B. A. (1990).
Adaptive selling: Conceptualization, measurement, and nomological validity.
\textit{Journal of Marketing Research, 27}(1), 61--69.

\bibitem{zohaib2018}
Zohaib, M. (2018).
Dynamic difficulty adjustment (DDA) in computer games: A review.
\textit{Advances in Human-Computer Interaction, 2018}, 1--12.

\end{thebibliography}

\end{document}
